%% 第24回編集委員会（物理第3巻）：組 波W6
%% p3-12.tex  波-12　光の進み方：反射・屈折・レンズ（★1〜4・3段にしない）
%% 中身は計画書 tools/phys3-plan.md の B「波-12」。公式の持ち主は A（A2・A7・A8・A9・A10・A33〜A37・A61）。
%%   反射・屈折の法則（波一般、A28）の持ち主は波-8 ①。この章は光に当てはめ、フェルマーの原理から 5-7 の最小で導き直す。
%% 数値はすべて Python で計算した：checks/物理第3巻/波-12/calc_p3_12.py（出力 calc_p3_12_out.txt）、ray_code.py（⑦の囲みのプログラム）
%% 記号（この章。第4部）：真空の光速 c（W-4。第1部の減衰係数・第2〜3部の波の速さとは別）、物質の中の光の速さ v＝c/n、屈折率 n、
%%   入射角・屈折角 θ₁・θ₂、臨界角 θ_c、焦点距離 f（この章では振動数 f を使わない）、物と像の距離 a・b、曲率半径 R、
%%   光線の高さ y と傾き θ（7-3 と同じ）、行列の成分 A・B・C・D、プリズムの頂角 α・偏角 δ、光束 Φ・光度 I・照度 E、度数（ディオプトリー）D は文字にしない。
%% 波-13 との約束：③で n＝c/v と屈折の法則、⑧で「プリズムでは紫の光ほどよく曲がる（分散）」を書く（まとめ役の連絡・組どうしの確かめ）。
%%   可視光の範囲 約 380〜780 nm は熱-13・波-13 と同じ。
%% 数学の本が引いている所：7-3 の光線の行列（すき間 [[1,d],[0,1]]・レンズ [[1,0],[-1/f,1]]、「（レンズ）×（すき間）」の順。7-3 の章末問題 d=0.1・f=0.05）、
%%   7-1 の「カメラのレンズと光の道すじ（物理 波-12）」、5-17 の sinθ≈θ（光の屈折）、5-7 の最大・最小。
%% 出典（第24回の追記（11月5日）により、書くときに確かめた。WebSearch で2つ以上の資料を照らし合わせたものは「照合」と書いた）
%%   ・真空の光速 c＝299 792 458 m/s（定義値、exact）：BIPM の SI 基本単位のページ（https://www.bipm.org/en/si-base-units/metre）と NIST（https://physics.nist.gov/cgi-bin/cuu/Value?c）を curl で読んだ。
%%       1983年の第17回国際度量衡総会で、メートルを光の進む距離で決めた（poynton.ca・arXiv:1606.02128 で照合）。
%%   ・レーマー（1676年、木星の衛星イオの食の時刻のずれ。地球の軌道を横切るのに約22分）、フィゾー（1849年、歯の数720の歯車、
%%       約 8.6 km 先の鏡、秒速約31万 km）：Linda Hall Library の Scientist of the Day「Hippolyte Fizeau」（検索の抜粋）、scienceready.com.au で照合。
%%       22分は当時の見積もりで、今の値（約16.6分）より長いことを本文に書いた。
%%   ・水の屈折率（修3-5、P5-22）：IAPWS R9-97（Release on the Refractive Index of Ordinary Water Substance, 1997。Harvey ほか J. Phys. Chem. Ref. Data 27, 761 (1998)）の式で計算した。
%%       iapws.org と NIST の CRC 用の表（tsapps.nist.gov）は作業環境から届かなかった（プロキシが 403。11月5日）。そこで、R9-97 を実装した Python の chemicals 1.5.2
%%       （PyPI）の係数で計算し、その実装の検算値（25 ℃・997.047 kg/m³・589.3 nm で 1.332858）と合うことを確かめた。水の密度は NIST WebBook（IAPWS-95）。
%%       20 ℃（真空に対して）：589.3 nm 1.33335、656.3 nm 1.33151、405 nm 1.34309。空気から水へ入る光で効く、空気に対する値（空気 n≈1.000272 でわる）は
%%       1.33299・1.33115・1.34273 で、本文の 1.333・1.331・1.343 と合う（checks/物理第3巻/確認/fix5_iapws_water_n.py）。下の二次の資料の値とも合う。
%%   ・屈折率（黄色の光 約 589 nm）：水 1.333（Tilton & Taylor の値を引く engineersedge.com の表 20 ℃で 1.33336、Dorsey の表を引く scubageek.com で 1.3330）、
%%       ダイヤモンド 2.417（本文は 2.42。gemsociety.org の宝石の屈折率の表、カリフォルニア大学バークレー校 EPS2 の講義資料）、
%%       ガラス（SCHOTT N-BK7 のデータシート nd＝1.51680。本文は「約1.5」）、空気 1.0003（schoolphysics.co.uk・hypertextbook.com）。
%%   ・水の屈折率の波長による変化（20 ℃）：656 nm で 1.331、405 nm 付近で 1.343（engineersedge.com の 20 ℃の表で 404 nm 1.34315・633 nm 1.33211、
%%       Dorsey の表を引く scubageek.com で 404.7 nm 1.3435 ほか。2つの表の差は 0.001 以内）。
%%   ・光ファイバ（ステップ形のマルチモード）のコア 1.48・クラッド 1.46、NA 0.24・受け入れ角 約14°：hansrajcollege.ac.in の講義資料、
%%       firgelliauto.com の計算例で照合。
%%   ・JIS Z 9110:2010（照明基準総則）の事務所の事務室 750 lx（推奨照度・維持照度）：九州電力の資料、パナソニックの照明設計の資料、
%%       環境省の資料（事務所 500〜1000 lx の範囲）で照合。JIS の原文は閲覧していない。
%%   ・LED 電球の60形相当 810 lm：日立の2012年の発表資料（日本電球工業会のガイドラインによる）、東芝の製品の表示で照合。
%%   ・晴れた日の、太陽に向けた面の照度 約10万 lx（直射）：パナソニックの昼光シミュレーションの資料、照明学会誌の論文（大気圏外 約14万 lx）で照合。
%%   ・明視の距離 25 cm：虫めがねの倍率 25/f の約束（高校物理の教科書の書き方）。
%%   ・ミツバチが空の偏光で太陽の向きを知ること：von Frisch, Experientia (1949)。昆虫の脳の「中心複合体」が方角を受け持つが、
%%       どのように進む向きを決めているかはまだはっきりしない：PMC5201003 ほかの総説、J. Exp. Biol. 223 jeb228254 で照合。
%%   ・空気の屈折率 15 ℃・101.325 kPa で 532 nm 1.00027821、1064 nm 1.00027398（LIGO の資料 G2401406「Dispersion in Air」。Ciddor (1996) の式、
%%       NIST の Ciddor の計算ページを指す）。589 nm は約 1.000277。屈折率−1 が密度に比例する近似（グラッドストン・デールの関係）で 50 ℃の値を見積もった。
%%   ・フーコー（1850年、回る鏡）が水の中の光が空気より遅いことを示した：MacTutor（mathshistory.st-andrews.ac.uk の Foucault の伝記）、EPFL graphsearch で照合。
%%       フェルマー（1607〜1665年）、1662年ごろに屈折を最短時間で導いた：MacTutor の Fermat の伝記などの一般の科学史による（本文は「ごろ」）。
%%   ・〔要確認〕なし。
\makeatletter
% 光線（矢印を途中につける）  \rxmray[色]{(始め)}{(終わり)}
\newcommand\rxmray[3][bkRed]{\draw[#1, line width=1pt, -{Stealth[length=2.2mm]}] #2 -- ($#2!0.56!#3$);\draw[#1, line width=1pt] ($#2!0.5!#3$) -- #3;}
% 凸レンズ（たて）  \rxmlens{x}{半分の高さ}
\newcommand\rxmlens[2]{\draw[bkBlue!70!black, line width=0.9pt, fill=bkBlue!12] (#1,-#2) .. controls (#1+2.6,-#2*0.4) and (#1+2.6,#2*0.4) .. (#1,#2) .. controls (#1-2.6,#2*0.4) and (#1-2.6,-#2*0.4) .. (#1,-#2);}
% 凹レンズ（たて）  \rxmclens{x}{半分の高さ}
\newcommand\rxmclens[2]{\draw[bkBlue!70!black, line width=0.9pt, fill=bkBlue!12] (#1-1.8,-#2) -- (#1+1.8,-#2) .. controls (#1+0.4,-#2*0.4) and (#1+0.4,#2*0.4) .. (#1+1.8,#2) -- (#1-1.8,#2) .. controls (#1-0.4,#2*0.4) and (#1-0.4,-#2*0.4) .. (#1-1.8,-#2);}
% 矢印の形の物（たてに立つ）  \rxmobj{(根もと)}{高さ}{色}
\newcommand\rxmobj[3]{\draw[#3, line width=1.6pt, -{Stealth[length=2.6mm, width=2.2mm]}] #1 -- ++(0,#2);}
% 目  \rxmeye{(位置)}{向きの角}
\newcommand\rxmeye[2]{\begin{scope}[shift={#1}, rotate=#2]\draw[bkInk, line width=0.7pt, fill=white] (0,0) .. controls (-2,1.6) and (-4.4,1.4) .. (-5,0) .. controls (-4.4,-1.4) and (-2,-1.6) .. (0,0);\fill[bkInk] (-3.4,0) circle (0.7);\end{scope}}
\makeatother

\phchapter{12}{光の進み方：\\反射・屈折・レンズ}{水に入れたストローは、なぜ折れて見える？}{\bklevelrange{1}{4}}{光の進み方：反射・屈折・レンズ}
\bkmichishirube{4-3,4-6,5-7,5-17,7-3,波-8}{波-13,電-15,原-2}

\begin{goal}
\begin{itemize}[label={\color{bkTheme}●}, leftmargin=1.6\zw, itemsep=0pt, topsep=1pt]
\item \bklevel{1}\enspace 光の直進・反射・屈折を、身近な例と作図で説明できる。
\item \bklevel{3}\enspace 屈折の法則で、屈折角・臨界角・見かけの深さを求められる。
\item \bklevel{3}\enspace レンズの式で、像の位置と倍率を求め、光学機器の倍率を見積もれる。
\item \bklevel{3}\enspace 屈折率が色で変わること（分散）から、プリズムと虹の色を説明できる。
\item \bklevel{4}\enspace フェルマーの原理から反射・屈折の法則を導き、光線の行列でレンズの組を計算できる。
\end{itemize}
\end{goal}

\phkoushiki{%
\phk{反射の法則（入射角＝反射角）}{中学理科}{ph:波-12:hansha}
\phk{光の屈折、全反射}{中学理科}{ph:波-12:kussetsu}
\phk{凸レンズの像（焦点、実像・虚像）}{中学理科}{ph:波-12:zou}
\phk{屈折率 $n=c/v$、屈折の法則、臨界角}{高校物理}{ph:波-12:snell}
\phk{レンズの式 $\tfrac1a+\tfrac1b=\tfrac1f$、倍率 $\tfrac ba$}{高校物理}{ph:波-12:lens}
\phk{球面鏡の式 $\tfrac1a+\tfrac1b=\tfrac2R$}{高校物理}{ph:波-12:kyumen}
\phk{光の分散・スペクトル・虹}{高校物理}{ph:波-12:bunsan}
}
%% 整合（U-5・K-8）：小学校理科の「わかること」の行（光の直進と反射）をけずった（本文①②はそのまま）。
%%   球面鏡は、高等学校学習指導要領（平成30年告示）解説 理科編「物理」(2)(ウ)㋐ に「凹面鏡や単一レンズの焦点と
%%   光の進路の規則性を扱う」とあるので、校種を「高校物理」にした（文部科学省 https://www.mext.go.jp/content/20250311-mxt_kyoiku02-100002620_05.pdf）。

\begin{hatena}
\noindent\begin{minipage}[c]{0.58\linewidth}\raggedright
コップに水を入れて、ストローをななめにさしてみましょう。横から見ると、ストローは水面の所で折れ曲がって見えます。水の中の部分が、浅い所へ浮き上がったように見えるのです。

ストローを抜いて見ると、まっすぐです。水が、ストローを曲げたのではありません。では、何が曲がっているのでしょう。
\end{minipage}\hfill
\begin{minipage}[c]{0.38\linewidth}\centering
\begin{tikzpicture}[x=1mm, y=1mm]
  \fill[bkBlue!15] (0,0) rectangle (26,20);
  \draw[bkInk, line width=0.8pt] (0,30) -- (0,0) -- (26,0) -- (26,30);
  \draw[bkBlue!70!black, line width=0.6pt] (0,20) -- (26,20);
  \draw[bkOrange!80!black, line width=2.4pt] (22,34) -- (14.5,20);
  \draw[bkOrange!80!black, line width=2.4pt] (14.5,20) -- (4,5.3);
  \draw[bkGray, line width=0.6pt, dashed] (14.5,20) -- (4,0.4);
  \node[font=\sffamily\footnotesize, anchor=west] at (27,20) {水面};
  \node[font=\sffamily\footnotesize, anchor=north, text=bkGray, align=center] at (13,-1) {破線：本当の\\ストロー};
\end{tikzpicture}
\end{minipage}
\end{hatena}

\noindent この章から、第4部「光」に入ります。第2部（波-5〜8）では、波が反射し、屈折し、重なることを学びました。光も波です（波-13）。でも、光の波長は1万分の数 mm ととても短いので、ふだんの大きさの物を相手にするときは、光を「まっすぐ進む線（\textbf{光線}）」として扱うと、ほとんどのことが説明できます。\bkidx{こうせん}{光線}この章では、光線の考えで、鏡・水・レンズ・プリズム・虹・カメラ・望遠鏡を読みときます。

\begin{yattemiyou}[コップの水とコインで、光の曲がりを見よう]
\textbf{用意するもの}\par
不透明なカップ、コイン、水さし、虫めがね、白い紙\par\vspace{2pt}
\noindent\begin{minipage}[c]{0.6\linewidth}
\begin{enumerate}[label={\bkcircled[bkGreen]{\arabic*}}, leftmargin=2.4\zw, itemsep=1pt, topsep=0pt]
\item カップの底にコインを置き、コインがカップのふちにかくれて、ちょうど見えなくなる所まで、顔を下げる。
\item 頭を動かさずに、友だちに、カップへ静かに水を入れてもらう。
\item 虫めがねで、窓の外の景色を白い紙に映す。くっきり映るときの、虫めがねと紙の距離をはかる。
\end{enumerate}
\end{minipage}\hfill
\begin{minipage}[c]{0.37\linewidth}\centering
\begin{tikzpicture}[x=0.9mm, y=0.9mm]
  \fill[bkBlue!15] (0,0) rectangle (24,14);
  \draw[bkInk, line width=0.8pt] (0,22) -- (0,0) -- (24,0) -- (24,22);
  \fill[bkYellow!80!black] (5,0) rectangle (11,1.2);
  \draw[bkRed, line width=0.9pt] (8,1.2) -- (17,14) -- (36,30);
  \draw[bkGray, line width=0.6pt, dashed] (24,22) -- (36,30);
  \rxmeye{(40,33)}{215}
  \node[font=\sffamily\footnotesize, anchor=north] at (8,-0.8) {コイン};
\end{tikzpicture}
\end{minipage}
\end{yattemiyou}

\noindent やってみようの答え：水を入れると、見えなかったコインが、浮き上がって見えてきます。コインから出た光が、水面で曲がって目に届くようになるからです（図の実線）。虫めがねでは、紙までの距離がある値（ふつうの虫めがねで\qtyrange{10}{20}{cm}くらい）のとき、景色が上下・左右さかさまにくっきり映ります。この距離が、虫めがねの焦点距離です（⑥）。

\begin{omoidasou}
\bkrecall{4-3}{相似な三角形では、対応する辺の比が等しい。}
\bkrecall{4-6}{$\sin\theta$ は斜辺に対する、向かいの辺の比。電卓で $\sin^{-1}$ を使うと角が出る。}
\bkrecall{5-7}{関数がいちばん小さくなる所では、導関数が0。}
\bkrecall{5-17}{小さな角（ラジアン）では $\sin\theta\approx\tan\theta\approx\theta$。}
\bkrecall{7-3}{光線を（高さ $y$、傾き $\theta$）の組で表す。すき間（左）とレンズ（右）の行列は
\[\begin{pmatrix}1&d\\0&1\end{pmatrix},\qquad\begin{pmatrix}1&0\\-1/f&1\end{pmatrix}\]}
\bkrecall{波-8}{ホイヘンスの原理から、境で波が曲がる：$\sin\theta_1/\sin\theta_2=v_1/v_2$。境で振動数は変わらない。}
\end{omoidasou}

\noindent この章の記号の注意：第4部では、$c$ を\textbf{真空の光の速さ}に使います（第1部では減衰係数、第2・3部では波の速さに使いました）。焦点距離を $f$ と書くので、この章では振動数を $f$ と書きません。

%% ======================== 物理の道具 ========================
\bksection{doguu}{物理の道具}

\bksubsection{1}{光はまっすぐ進む}
\label{ph:波-12:sec1}
\label{ph:波-12:chokushin}日なたにできる影の形は、物の形とそっくりです。木もれ日が細い線になって見えることもあります。これらは、同じ物の中（空気の中など）では、\textbf{光がまっすぐ進む}ことを示しています。この性質を、\textbf{光の直進}といいます。\bkidx{ひかりのちょくしん}{光の直進}

\phboxhead{ピンホールカメラ：相似で像の大きさを求める}
箱の前に小さな穴をあけ、反対側に薄い紙（スクリーン）をはると、外の景色が上下さかさまに映ります（\textbf{ピンホールカメラ}）。物の上のはしから出た光は、穴を通ってまっすぐ進み、スクリーンの下のほうに届くからです。

\noindent\begin{minipage}[c]{0.56\linewidth}\centering
\begin{tikzpicture}[x=1mm, y=1mm]
  \rxmobj{(0,0)}{20}{bkGreen!60!black}
  \fill[bkGray!40] (40,-14) rectangle (41.2,-1) (40,1) rectangle (41.2,24);
  \draw[bkInk, line width=0.8pt] (41.2,-14) -- (60,-14) -- (60,24) -- (41.2,24);
  \draw[bkOrange!80!black, line width=1.6pt, -{Stealth[length=2.2mm]}] (60,0) -- (60,-10);
  \draw[bkRed, line width=0.8pt] (0,20) -- (60,-10);
  \draw[bkRed, line width=0.8pt] (0,0) -- (60,0);
  \draw[bkInk, line width=0.5pt, {Stealth[length=2pt]}-{Stealth[length=2pt]}] (0,-17) -- (40.6,-17);
  \draw[bkInk, line width=0.5pt, {Stealth[length=2pt]}-{Stealth[length=2pt]}] (40.6,-17) -- (60,-17);
  \node[font=\sffamily\footnotesize, anchor=north] at (20,-17.5) {物までの距離};
  \node[font=\sffamily\footnotesize, anchor=north] at (50.3,-17.5) {箱の深さ};
  \node[font=\sffamily\footnotesize, anchor=east] at (-1,10) {物};
  \node[font=\sffamily\footnotesize, anchor=west] at (61,-5) {像};
  \node[font=\sffamily\footnotesize, anchor=south] at (40.6,24.5) {穴};
\end{tikzpicture}
\end{minipage}\hfill
\begin{minipage}[c]{0.4\linewidth}\raggedright
物と像は、穴をはさんで向かい合う、相似な三角形（4-3）です。
\[\frac{\text{像の大きさ}}{\text{物の大きさ}}=\frac{\text{箱の深さ}}{\text{物までの距離}}\]
\qty{10}{m}先の高さ\qty{5}{m}の木を、深さ\qty{20}{cm}の箱で見ると、像は $5\times0.2\div10=\qty{0.1}{m}$、\qty{10}{cm}です。
\end{minipage}

\phboxhead{光の速さ}
光はとても速く進みます。昔の人は、光は一瞬で届くと考えていました。光の速さがはじめて見積もられたのは、1676年です。デンマークのレーマーは、木星の衛星イオが木星のかげにかくれる時刻が、地球が木星から遠い季節には、予定よりおくれることに気づきました。光が地球の軌道の分だけ長い道のりを進むからだと考え、地球の軌道を横切るのに約22分かかると見積もりました（今の値は約16.6分）。地上で光の速さをはかったのは、1849年のフィゾーです。速く回る歯車のすき間から光を出し、約\qty{8.6}{km}先の鏡ではね返ってきた光が、次の歯にさえぎられるかどうかで往復の時間をはかり、秒速約31万 km を得ました。

今では、真空の中の光の速さ（\textbf{光速}）は、
\pheq{c=\num{299792458}\ \text{m/s}\quad（\text{約}\ \qty{3.00e8}{m/s}）}{m/s}
と\textbf{決められています}（1983年から。この数で長さの単位の 1\,m を決めている）。\bkidx{こうそく}{光速}光は1秒に地球を約7.5周し、月まで約1.3秒、太陽からは約8分20秒（499秒）で届きます。

\begin{phmijika}[雷の光と音]
雷が\qty{1}{km}先に落ちたとき、光は $1000\div(3\times10^8)\approx\qty{3e-6}{s}$ で、ほぼ一瞬で届きます。音（\qty{343}{m/s}）は約\qty{2.9}{s}かかります（波-9）。光ってから音が聞こえるまでの秒数に、約\qty{340}{m}をかければ、雷までの距離がわかります。
\end{phmijika}

\Needspace{10\baselineskip}\bksubsection{2}{反射の法則と、鏡の像}
\label{ph:波-12:sec2}
\label{ph:波-12:hansha}光が鏡に当たると、はね返ります（\textbf{反射}）。鏡の面に垂直な線（\textbf{法線}）と、入ってくる光のなす角を\textbf{入射角}、はね返った光のなす角を\textbf{反射角}といいます。\bkidx{はんしゃのほうそく}{反射の法則}実験すると、いつも
\pheq{\text{入射角}=\text{反射角}}{°}
になります。これを\textbf{反射の法則}といいます。入ってくる光・法線・はね返る光は、同じ平面の上にあります。

\phboxhead{なぜ等しいのか：いちばん短い道}
光が点 A から出て、鏡の点 P ではね返り、点 B に届くとします。鏡の上のどの点 P を通る道が、実際の光の道でしょう。B を鏡に対して折り返した点を B$'$ とすると、P から B までの長さは、P から B$'$ までの長さと同じです（鏡に対して対称だから）。
\[\text{AP}+\text{PB}=\text{AP}+\text{PB}'\]
右辺は、A から P を通って B$'$ へ行く折れ線の長さなので、P が A と B$'$ を結ぶ直線の上にあるとき、いちばん短くなります。そのとき、図\ref{ph:波-12:fig-hansha}の角 ① と ③ は対頂角で等しく、② と ③ は折り返しで等しいので、入射角＝反射角です。\textbf{光は、いちばん短い道（同じ物の中では、いちばん短い時間の道）を通る}と考えると、反射の法則が出てくるのです。この考えを\textbf{フェルマーの原理}といいます（③で詳しく）。

\begin{bkfigure}
\begin{tikzpicture}[x=1mm, y=1mm]
  \fill[bkGray!25] (-6,-1.6) rectangle (66,0);
  \draw[bkInk, line width=0.9pt] (-6,0) -- (66,0);
  \node[font=\sffamily\footnotesize, anchor=north west] at (-6,-1.8) {鏡};
  \draw[bkGray, line width=0.5pt, dashed] (30,-24) -- (30,24);
  \fill (6,20) circle (0.9) (54,20) circle (0.9) (54,-20) circle (0.9) (30,0) circle (0.8);
  \rxmray{(6,20)}{(30,0)}
  \rxmray{(30,0)}{(54,20)}
  \draw[bkRed!60, line width=0.8pt, dashed] (30,0) -- (54,-20);
  \draw[bkGray, line width=0.4pt, dotted] (54,20) -- (54,-20);
  \node[font=\sffamily\footnotesize, anchor=south east] at (5.5,20.5) {A};
  \node[font=\sffamily\footnotesize, anchor=south west] at (54.5,20.5) {B};
  \node[font=\sffamily\footnotesize, anchor=north west] at (54.5,-20.5) {B$'$};
  \node[font=\sffamily\footnotesize, anchor=north] at (27.5,-2) {P};
  \draw[bkInk, line width=0.5pt] (30,7) arc (90:140.2:7);
  \draw[bkInk, line width=0.5pt] (30,7) arc (90:39.8:7);
  \draw[bkInk, line width=0.5pt] (30,-7) arc (270:320.2:7);
  \node[font=\sffamily\footnotesize] at (25.6,10.6) {①};
  \node[font=\sffamily\footnotesize] at (34.4,10.6) {②};
  \node[font=\sffamily\footnotesize] at (34.4,-10.6) {③};
  \node[font=\sffamily\footnotesize, anchor=south] at (30,24.5) {法線};
  \node[anchor=west, font=\small, text width=46mm, align=left] at (72,4) {①入射角、②反射角。\par ②と③は鏡に対して折り返し、①と③は対頂角なので、①＝②。\par AP＋PB＝AP＋PB$'$ は、P が直線 AB$'$ の上でいちばん短い。};
\end{tikzpicture}
\bkfigcaption{反射の法則を、いちばん短い道で考える。B$'$ は B を鏡に対して折り返した点（破線は鏡の後ろへのばした光の道）。}\label{ph:波-12:fig-hansha}
\end{bkfigure}

\phboxhead{平面鏡の像}
鏡の前に立つと、鏡の後ろに自分がいるように見えます。物の各点から出て鏡で反射した光は、どれも、鏡に対して折り返した点（図\ref{ph:波-12:fig-hansha}の B$'$）から出たように進むからです。だから、平面鏡の像は、\textbf{鏡に対して物と対称な位置に、同じ大きさで}できます。実際に光が集まっているわけではない、見かけの像です（⑥の虚像）。

\begin{phmijika}[全身を映す鏡の高さ]
身長の全部を映すのに要る鏡の高さは、身長のちょうど半分です。頭のてっぺんから出た光は、目と頭の高さのまん中の高さで鏡に反射して目に届き、つま先から出た光は、目とつま先のまん中の高さで反射します。2つのまん中の差は、身長のちょうど半分です。鏡から何 m 離れても変わりません（例題1）。
\end{phmijika}

\phboxhead{乱反射と再帰反射}
紙や壁の表面は、細かく見るとでこぼこしているので、光はいろいろな向きにはね返ります（\textbf{乱反射}）。\bkidx{らんはんしゃ}{乱反射}どの向きから見ても物が見えるのは、乱反射のおかげです。一方、2枚の鏡を直角に合わせると、どの向きから入った光も、2回反射して、来た向きへそのまま帰っていきます（図\ref{ph:波-12:fig-saiki}）。3枚の鏡を互いに直角に組んだ「角」を細かく並べたものが、自転車の反射板や道路の標識に使われています（\textbf{再帰反射}）。\bkidx{さいきはんしゃ}{再帰反射}車のライトの光が、ちょうど運転する人の方へ返るので、夜でも明るく見えるのです。

\begin{bkfigure}
\begin{tikzpicture}[x=1mm, y=1mm]
  \draw[bkInk, line width=1.2pt] (0,30) -- (0,0) -- (30,0);
  \fill[bkGray!25] (-1.4,0) rectangle (0,30) (0,-1.4) rectangle (30,0);
  \rxmray{(40,11)}{(18,0)}
  \draw[bkRed, line width=1pt] (18,0) -- (0,9);
  \draw[bkRed, line width=1pt, -{Stealth[length=2.2mm]}] (0,9) -- (40,29);
  \rxmray[bkBlue!70!black]{(40,24)}{(14,0)}
  \draw[bkBlue!70!black, line width=1pt] (14,0) -- (0,12.92);
  \draw[bkBlue!70!black, line width=1pt, -{Stealth[length=2.2mm]}] (0,12.92) -- (26,36.92);
  \node[anchor=west, font=\small, text width=72mm, align=left] at (48,16) {直角に合わせた2枚の鏡（横から見た図）。\par 2本の光は、どちらも2回反射すると、入ってきた向きと平行に、反対向きに帰る。};
\end{tikzpicture}
\bkfigcaption{再帰反射。1回めの反射と2回めの反射で、光の向きはあわせて \ang{180} 変わる。}\label{ph:波-12:fig-saiki}
\end{bkfigure}

\Needspace{10\baselineskip}\bksubsection{3}{屈折の法則：光は「早く着く道」を選ぶ}
\label{ph:波-12:sec3}
\label{ph:波-12:kussetsu}光が空気から水やガラスへななめに入ると、境の所で進む向きが変わります。これを\textbf{屈折}といいます。\bkidx{くっせつ}{屈折}法線と、屈折したあとの光のなす角を\textbf{屈折角}といいます。

\begin{bkfigure}
\begin{tikzpicture}[x=1mm, y=1mm]
  \fill[bkBlue!12] (0,0) rectangle (60,-26);
  \draw[bkBlue!70!black, line width=0.7pt] (0,0) -- (60,0);
  \draw[bkGray, line width=0.5pt, dashed] (30,26) -- (30,-26);
  \rxmray{(8,22)}{(30,0)}
  \rxmray{(30,0)}{(43.7,-22)}
  \draw[bkRed!50, line width=0.6pt, dashed] (30,0) -- (52,-22);
  \draw[bkInk, line width=0.5pt] (30,8) arc (90:135:8);
  \draw[bkInk, line width=0.5pt] (30,-8) arc (270:302:8);
  \node[font=\sffamily\footnotesize] at (26.6,11.6) {$\theta_1$};
  \node[font=\sffamily\footnotesize] at (33.2,-11.6) {$\theta_2$};
  \node[font=\sffamily\footnotesize, anchor=west] at (61,10) {空気（屈折率 $n_1$）};
  \node[font=\sffamily\footnotesize, anchor=west] at (61,-10) {水（屈折率 $n_2$）};
  \node[font=\sffamily\footnotesize, anchor=south] at (30,26.5) {法線};
  \node[font=\sffamily\footnotesize, anchor=north west, text=bkRed!70!black] at (50,-22.5) {まっすぐ進んだら};
\end{tikzpicture}
\bkfigcaption{空気から水へ入る光の屈折。光は、遅い物（水）へ入ると、法線に近づく向きに曲がる（$\theta_2<\theta_1$）。}\label{ph:波-12:fig-kussetsu}
\end{bkfigure}

\phboxhead{屈折率}
光は、水やガラスの中では、真空の中より遅く進みます。真空の中の光速 $c$ が、物の中の光の速さ $v$ の何倍かを、その物の\textbf{屈折率}といいます。\bkidx{くっせつりつ}{屈折率}
\pheq{n=\frac{c}{v}\qquad\text{（物の中の光の速さ}\ v=c/n\text{）}}{1 ＝ $\dfrac{\text{m/s}}{\text{m/s}}$}
\label{ph:波-12:snell}
\begin{bkcenter}
\begin{bktable}{colspec={Q[l,m,wd=30mm] *{5}{Q[c,m,wd=14.5mm]}}, rowsep=1.5pt, row{1}={bg=bkTheme!10, font=\sffamily\small}}
物 & 真空 & 空気 & 水 & ガラス & ダイヤモンド \\
屈折率 $n$ & 1（定義） & 1.0003 & 1.333 & 約1.5 & 2.42 \\
光の速さ〔$10^8$\,m/s〕 & 3.00 & 3.00 & 2.25 & 2.0 & 1.24 \\
\end{bktable}
\bktabcaption{いろいろな物の屈折率（黄色のナトリウムの光、波長 \qty{589}{nm}）。空気はほとんど1なので、この章では空気の屈折率を1とする。}\label{ph:波-12:tab-n}
\end{bkcenter}

\phboxhead{フェルマーの原理から屈折の法則を導く}
1662年ごろ、フェルマーは「光は、ある点から別の点へ、\textbf{かかる時間がいちばん短い道}を通る」と考えました（\textbf{フェルマーの原理}）。\bkidx{ふぇるまーのげんり}{フェルマーの原理}せいかくに言うと、光が通るのは、道を少し変えてもかかる時間がほとんど変わらない道（時間が極値になる道。停留する道ともいう）です。たいていは「いちばん短い」ですが、鏡の形によっては（強く曲がった凹面鏡など）、いちばん長い道になることもあります。同じ物の中では、時間がいちばん短い道は、長さがいちばん短い道（まっすぐ、②の反射）です。速さの違う2つの物にまたがるときは、どうなるでしょう。

\begin{phmijika}[ライフガードの走る道]
砂浜のライフガードが、海でおぼれている人を助けに行きます。砂の上を走るほうが、泳ぐより速いので、まっすぐ向かうより、\textbf{砂の上を少し長く走って、泳ぐ距離を短くする}ほうが、早く着きます。でも、走りすぎると遠回りです。いちばん早い道は、その間のどこかにあります（例題4）。光も、遅い水の中を進む距離を短くするように、境で向きを変えるのです。
\end{phmijika}

\begin{bkfigure}
\begin{tikzpicture}[x=1mm, y=1mm]
  \fill[bkYellow!20] (-4,0) rectangle (64,22);
  \fill[bkBlue!12] (-4,0) rectangle (64,-22);
  \draw[bkBlue!70!black, line width=0.7pt] (-4,0) -- (64,0);
  \fill (0,20) circle (0.9) (60,-20) circle (0.9);
  \node[font=\sffamily\footnotesize, anchor=south] at (0,20.5) {A};
  \node[font=\sffamily\footnotesize, anchor=north] at (60,-20.5) {B};
  \rxmray{(0,20)}{(43,0)}
  \rxmray{(43,0)}{(60,-20)}
  \draw[bkGray, line width=0.6pt, dashed] (0,20) -- (60,-20);
  \draw[bkGray, line width=0.5pt, dashed] (43,12) -- (43,-12);
  \draw[bkInk, line width=0.5pt, {Stealth[length=2pt]}-{Stealth[length=2pt]}] (0,-3) -- (43,-3);
  \node[font=\sffamily\footnotesize, anchor=north, fill=bkBlue!12, inner sep=1pt] at (21.5,-3.5) {$x$};
  \draw[bkInk, line width=0.5pt, {Stealth[length=2pt]}-{Stealth[length=2pt]}] (-3,0) -- (-3,20);
  \node[font=\sffamily\footnotesize, anchor=east] at (-3.5,10) {$h_1$};
  \draw[bkInk, line width=0.5pt, {Stealth[length=2pt]}-{Stealth[length=2pt]}] (63,0) -- (63,-20);
  \node[font=\sffamily\footnotesize, anchor=west] at (63.5,-10) {$h_2$};
  \draw[bkInk, line width=0.5pt] (43,7) arc (90:155.1:7);
  \draw[bkInk, line width=0.5pt] (43,-7) arc (270:310.4:7);
  \node[font=\sffamily\footnotesize] at (39.2,10) {$\theta_1$};
  \node[font=\sffamily\footnotesize] at (53,-5.5) {$\theta_2$};
  \node[anchor=west, font=\small, text width=46mm, align=left] at (70,8) {上：速さ $v_1$（砂浜・空気）\par 下：速さ $v_2$（海・水）\par A と B の横の距離 $L$};
  \node[anchor=west, font=\small, text width=46mm, align=left] at (70,-12) {破線：まっすぐの道。\par 遅い所の道が長いので、かえって時間がかかる。};
\end{tikzpicture}
\bkfigcaption{A から境の点（A から横に $x$）を通って B へ行く道。かかる時間がいちばん短くなる $x$ を探す。}\label{ph:波-12:fig-fermat}
\end{bkfigure}

\noindent A は境から高さ $h_1$、B は境から深さ $h_2$ の所にあり、横に $L$ 離れています。境の点を、A から横に $x$ の所にとると、かかる時間は
\[T(x)=\frac{\sqrt{h_1^{\,2}+x^2}}{v_1}+\frac{\sqrt{h_2^{\,2}+(L-x)^2}}{v_2}\]
です（距離÷速さ、長さは三平方の定理 4-4）。$T$ がいちばん小さい所では、$\dd T/\dd x=0$ です（5-7）。合成関数の微分（5-6）で、
\begin{bkhenkei}[0.5]
\frac{\dd T}{\dd x}=\frac{x}{v_1\sqrt{h_1^{\,2}+x^2}}-\frac{L-x}{v_2\sqrt{h_2^{\,2}+(L-x)^2}} & $\sqrt{u}$ の微分は $\dfrac{u'}{2\sqrt u}$ \\
\frac{x}{\sqrt{h_1^{\,2}+x^2}}=\sin\theta_1 & 図の直角三角形（4-6） \\
\frac{L-x}{\sqrt{h_2^{\,2}+(L-x)^2}}=\sin\theta_2 & 下の直角三角形も同じ \\
\frac{\sin\theta_1}{v_1}-\frac{\sin\theta_2}{v_2}=0 & $\dd T/\dd x=0$ とおく \\
\frac{n_1\sin\theta_1}{c}=\frac{n_2\sin\theta_2}{c} & $v=c/n$ を入れる \\
\end{bkhenkei}
両辺に $c$ をかけて、
\pheq{n_1\sin\theta_1=n_2\sin\theta_2}{1 ＝ 1}
が出ます。これを\textbf{屈折の法則}（スネルの法則）といいます。\bkidx{くっせつのほうそく}{屈折の法則}2回微分すると、どの $x$ でも正（下に凸、5-7）なので、ここが本当に「いちばん短い」所です。波-8 ①では、同じ法則を、ホイヘンスの原理（波面の各点から出る小さな波）から $\sin\theta_1/\sin\theta_2=v_1/v_2$ として導きました。$v=c/n$ を入れれば、まったく同じ式です。2つの違う考え方が、同じ法則に行き着くのです。

\begin{phyomi}[$n_1\sin\theta_1=n_2\sin\theta_2$]
「屈折率×（法線からの角の $\sin$）は、境の両側で同じ」。屈折率の大きい物（光が遅い物）では、$\sin\theta$ が小さく、光は法線に近づきます。光の進む向きを逆にしても同じ式なので、光の道は逆向きにもたどれます（\textbf{光の逆進性}）。
\end{phyomi}

\begin{phbaisuru}{入射角 $\theta_1$}
屈折角は2倍にはなりません。$\sin\theta_2=\sin\theta_1/n$ なので、入射角が小さいうちは屈折角もほぼ比例して増えますが、入射角が大きくなるにつれて、増え方がにぶります。空気から水へ入るときは、入射角を\ang{90}近くまで大きくしても、屈折角は約\ang{48.6}より大きくなりません（図\ref{ph:波-12:fig-kakudo}）。この\ang{48.6}が、④の臨界角です。
\par\smallskip\noindent\begin{tikzpicture}
\begin{axis}[bkplot, width=100mm, height=50mm, xmin=0, xmax=95, ymin=0, ymax=62, xtick={15,30,45,60,75,90}, ytick={15,30,41.8,48.6}, yticklabel style={/pgf/number format/fixed}, xlabel={入射角 $\theta_1$（°）}, ylabel={屈折角 $\theta_2$（°）}]
  \addplot[bkGray, line width=0.6pt, dashed, domain=0:62] {x};
  \addplot[bkcurve, domain=0:90, samples=60] {asin(sin(x)/1.333)};
  \addplot[bkcurve2, dash pattern=on 4pt off 2pt, domain=0:90, samples=60] {asin(sin(x)/1.5)};
  \node[font=\sffamily\footnotesize, anchor=south east, text=bkBlue] at (axis cs:90,55) {水（実線）};
  \node[font=\sffamily\footnotesize, anchor=north east, text=bkOrange!80!black] at (axis cs:90,33) {ガラス（破線）};
  \node[font=\sffamily\footnotesize, anchor=south east, text=bkGray] at (axis cs:53,56.5) {曲がらなかったら};
\end{axis}
\end{tikzpicture}
\bkfigcaption{入射角と屈折角（空気から水・ガラスへ）。屈折角は臨界角（水 約\ang{48.6}、ガラス 約\ang{41.8}）をこえない。}\label{ph:波-12:fig-kakudo}
\end{phbaisuru}

\phboxhead{はてな？の答え：見かけの深さ}
\label{ph:波-12:hatena}水の中のストローの先（深さ $h$）から出た光は、水面で、法線から遠ざかる向きに曲がって目に届きます（水から空気へ、$\theta_1$ が水の中の角）。目は、光がまっすぐ来たと思うので、光を逆にのばした所、つまり、本当より浅い所（深さ $h'$）にストローの先があると感じます。

\noindent\begin{minipage}[c]{0.46\linewidth}\centering
\begin{tikzpicture}[x=1mm, y=1mm]
  \fill[bkBlue!12] (-14,0) rectangle (40,-32);
  \draw[bkBlue!70!black, line width=0.7pt] (-14,0) -- (40,0);
  \fill (0,-30) circle (1);
  \fill[bkGray] (0,-21) circle (1);
  \rxmray{(0,-30)}{(12,0)}
  \rxmray{(12,0)}{(30,31.5)}
  \draw[bkGray, line width=0.6pt, dashed] (12,0) -- (0,-21);
  \draw[bkGray, line width=0.5pt, dashed] (0,6) -- (0,-32);
  \draw[bkInk, line width=0.5pt, {Stealth[length=2pt]}-{Stealth[length=2pt]}] (-11,0) -- (-11,-30);
  \node[font=\sffamily\footnotesize, anchor=east] at (-11.5,-15) {$h$};
  \draw[bkInk, line width=0.5pt, {Stealth[length=2pt]}-{Stealth[length=2pt]}] (-4,0) -- (-4,-21);
  \node[font=\sffamily\footnotesize, anchor=east] at (-4.5,-10.5) {$h'$};
  \draw[bkInk, line width=0.5pt, {Stealth[length=2pt]}-{Stealth[length=2pt]}] (0,2) -- (12,2);
  \node[font=\sffamily\footnotesize, anchor=south] at (6,2.3) {$s$};
  \node[font=\sffamily\footnotesize, anchor=south] at (30,32) {目へ};
  \node[font=\sffamily\footnotesize, anchor=west] at (13,-29) {本当の先};
  \node[font=\sffamily\footnotesize, anchor=west, text=bkGray] at (13,-20) {見える所};
\end{tikzpicture}
\end{minipage}\hfill
\begin{minipage}[c]{0.5\linewidth}\raggedright
水面で光が出る点までの横の距離 $s$ は、
\[s=h\tan\theta_1=h'\tan\theta_2\]
です（$\theta_1$ は水の中、$\theta_2$ は空気の中の、法線からの角）。真上に近い向きから見るときは、角が小さいので $\tan\theta\approx\sin\theta$（5-17）。屈折の法則 $n\sin\theta_1=\sin\theta_2$ を使うと、
\[\frac{h'}{h}=\frac{\tan\theta_1}{\tan\theta_2}\approx\frac{\sin\theta_1}{\sin\theta_2}=\frac1n\]
\end{minipage}
\par\vspace{2pt}
\Needspace{12\baselineskip}\begin{phflow}[屈折の法則から]
\phflowrow{$n_1\sin\theta_1=n_2\sin\theta_2$}{真上に近い向きから見る\\（小さな角、5-17）}{見かけの深さ $h'=\dfrac{h}{n}$}
\end{phflow}
水（$n=1.333$）なら、見かけの深さは本当の $1/1.333\approx0.75$、約4分の3です。深さ\qty{1}{m}のプールは、真上から見ると約\qty{75}{cm}に見えます。ストローの水の中の部分は、どの点も浅い所へ浮き上がって見えるので、水面の所で折れて見えるのです。これがはてな？の答えです。曲がっているのは、ストローではなく、\textbf{光の道}でした。ななめから見るほど、浮き上がり方は大きくなります（空気の側の角が\ang{40}で約0.66倍、\ang{60}で約0.49倍）。

\Needspace{10\baselineskip}\bksubsection{4}{全反射と光ファイバ}
\label{ph:波-12:sec4}
光が、屈折率の大きい物から小さい物へ（水から空気へ）出るときは、法線から遠ざかる向きに曲がります。入射角を大きくしていくと、屈折角が先に\ang{90}になり、それより大きな入射角では、光は空気へ出られず、境で全部はね返ります。これを\textbf{全反射}といいます。\bkidx{ぜんはんしゃ}{全反射}屈折角がちょうど\ang{90}になる入射角を\textbf{臨界角} $\theta_{\mathrm c}$ といいます。\bkidx{りんかいかく}{臨界角}屈折の法則で $\theta_2=\ang{90}$（$\sin\theta_2=1$）とおくと、
\pheq{\sin\theta_{\mathrm c}=\frac{n_2}{n_1}\qquad\text{（空気へ出るときは}\ \sin\theta_{\mathrm c}=\frac1n\text{）}}{1}
\label{ph:波-12:rinkai}
\Needspace{8\baselineskip}\begin{phflow}[臨界角]
\phflowrow{$n_1\sin\theta_1=n_2\sin\theta_2$}{$n_1>n_2$、\\屈折角 $\theta_2=\ang{90}$}{$\sin\theta_{\mathrm c}=\dfrac{n_2}{n_1}$（空気へは $\dfrac1n$）}
\end{phflow}

\begin{bkfigure}
\begin{tikzpicture}[x=1mm, y=1mm]
  \fill[bkBlue!12] (-30,0) rectangle (50,-24);
  \draw[bkBlue!70!black, line width=0.7pt] (-30,0) -- (50,0);
  \fill[bkYellow!80!black] (0,-22) circle (1.4);
  \draw[bkRed, line width=0.9pt] (0,-22) -- (-8,0) -- (-15.2,14);
  \draw[bkOrange!80!black, line width=0.9pt] (0,-22) -- (-16,0) -- (-27.4,9);
  \draw[bkPurple, line width=1.2pt] (0,-22) -- (24.9,0) -- (40,0.3);
  \draw[bkBlue!70!black, line width=0.9pt] (0,-22) -- (30,0) -- (50,-14.7);
  \node[font=\sffamily\footnotesize, anchor=north west] at (1.6,-21) {光を出す点};
  \node[font=\sffamily\footnotesize, anchor=south, text=bkPurple] at (34,0.8) {臨界角：水面にそう};
  \node[font=\sffamily\footnotesize, anchor=north west, text=bkBlue!70!black] at (36,-16) {全反射};
  \node[font=\sffamily\footnotesize, anchor=south] at (-15.2,14.5) {屈折して出る};
  \node[font=\sffamily\footnotesize, anchor=west] at (51,4) {空気};
  \node[font=\sffamily\footnotesize, anchor=west] at (51,-20) {水};
\end{tikzpicture}
\bkfigcaption{水の中の点から出た光。入射角が小さい光は屈折して空気へ出る。臨界角（水は約\ang{48.6}）で水面にそって進み、それより大きいと全部はね返る（全反射）。}\label{ph:波-12:fig-zenhansha}
\end{bkfigure}

\noindent 水の臨界角は $\sin\theta_{\mathrm c}=1/1.333$ より約\ang{48.6}、ガラス（1.5）は約\ang{41.8}、ダイヤモンド（2.42）は約\ang{24.4}です（例題5）。ダイヤモンドは臨界角がとても小さいので、中に入った光が何度も全反射してから上へ出てきます。ダイヤモンドがきらきら輝くのは、このためです。

\begin{phmijika}[水の中から見上げた空]
プールにもぐって水面を見上げると、空は、頭の上の円い窓の中にだけ見えます。窓の外側は、水の底が全反射で映って見えます。空から来るどんな向きの光も、水に入ると、法線から\ang{48.6}以内に曲げられるからです。深さ\qty{2}{m}なら、窓の半径は $2\times\tan\ang{48.6}\approx\qty{2.3}{m}$ です。
\end{phmijika}

\begin{motto}[4]
\textbf{逃げ水：空気の中の全反射}\par
夏の晴れた日、遠くの道路に水たまりがあるように見え、近づくと消えてしまうことがあります（\textbf{逃げ水}）。\bkidx{にげみず}{逃げ水}熱い路面の近くの空気は温まって密度が小さく、屈折率が少しだけ小さくなっています。空気の屈折率から1を引いた値（\qty{15}{\celsius}で約 $2.8\times10^{-4}$）は、密度にほぼ比例するので、路面近くが\qty{50}{\celsius}なら、絶対温度の比 $288\div323$ をかけて約 $2.5\times10^{-4}$ で、差は約 $3\times10^{-5}$ です。

\noindent\begin{minipage}[c]{0.5\linewidth}\centering
\begin{tikzpicture}[x=1mm, y=1mm]
  \fill[bkGray!40] (0,-1.2) rectangle (62,0);
  \fill[bkOrange!15] (0,0) rectangle (62,4);
  \draw[bkRed, line width=1pt] (2,14) .. controls (30,0.2) and (34,0.2) .. (60,9);
  \draw[bkGray, line width=0.6pt, dashed] (60,9) -- (30,-1.2);
  \draw[bkInk, line width=0.6pt, fill=white] (60.5,9) circle (1.2);
  \node[font=\sffamily\footnotesize, anchor=south west] at (2,14) {空の光};
  \node[font=\sffamily\footnotesize, anchor=north] at (31,-2.2) {熱い路面（その上の空気は屈折率が小さい）};
  \node[font=\sffamily\footnotesize, anchor=south east] at (59,10.4) {目};
\end{tikzpicture}
\end{minipage}\hfill
\begin{minipage}[c]{0.46\linewidth}\raggedright\small
空から路面すれすれに来た光は、屈折率の小さい層で少しずつ上へ曲げられ、全反射のようにはね上がって目に届きます。目は、光を逆にのばした路面の上（破線）に、空が映った「水面」を見ます。
\end{minipage}

\par\smallskip
④の臨界角の式で $\cos\alpha=n_2/n_1$（$\alpha$ は路面となす角）とおくと、
\[\alpha\approx\sqrt{2\times3\times10^{-5}}\approx\qty{0.0077}{rad}\quad(\text{約}\,\ang{0.44})\]
です。路面とこれより小さな角をなす光だけがはね上がるので、目の高さ\qty{1.5}{m}なら、逃げ水は約 $1.5\div0.0077\approx\qty{190}{m}$ より先に見えます。近づくと見る角が大きくなるので、逃げ水は先へ逃げていくのです。
\end{motto}

\phboxhead{光ファイバ}
\textbf{光ファイバ}は、髪の毛ほどの細いガラスの線で、光を遠くまで運びます。\bkidx{ひかりふぁいば}{光ファイバ}まん中（\textbf{コア}、屈折率 $n_1$）を、屈折率が少し小さい外側（\textbf{クラッド}、$n_2$）が包んでいます。コアの中の光は、クラッドとの境で全反射を繰り返して、曲がった線の中でも外へもれずに進みます。

\begin{bkfigure}
\begin{tikzpicture}[x=0.8mm, y=1mm]
  \fill[bkBlue!8] (0,-10) rectangle (100,10);
  \fill[bkBlue!20] (0,-5) rectangle (100,5);
  \draw[bkBlue!70!black, line width=0.6pt] (0,-10) -- (100,-10) (0,10) -- (100,10);
  \draw[bkBlue!70!black, line width=0.4pt, dashed] (0,-5) -- (100,-5) (0,5) -- (100,5);
  \draw[bkGray, line width=0.4pt, dashed] (-26,0) -- (0,0);
  \draw[bkRed, line width=1pt] (-20,-4.9) -- (0,0) -- (21.3,5) -- (63.9,-5) -- (100,3.5);
  \draw[bkInk, line width=0.5pt] (-8,0) arc (180:193.8:8);
  \node[font=\sffamily\footnotesize, anchor=south] at (-12,0.6) {$\theta_{\mathrm a}$};
  \node[font=\sffamily\footnotesize, anchor=west] at (101,0) {コア $n_1$};
  \node[font=\sffamily\footnotesize, anchor=west] at (101,7.5) {クラッド $n_2$};
  \node[font=\sffamily\footnotesize, anchor=north] at (21.3,-11) {全反射を繰り返して進む};
\end{tikzpicture}
\bkfigcaption{光ファイバ（たての向きを大きくかいた）。はしに入る光の角 $\theta_{\mathrm a}$ が受け入れ角より小さければ、コアとクラッドの境で全反射を続ける。}\label{ph:波-12:fig-fiber}
\end{bkfigure}

\noindent はしから入る光の角 $\theta_{\mathrm a}$ が大きすぎると、コアとクラッドの境での入射角が臨界角より小さくなり、光はクラッドへもれます。全反射が続くいちばん大きな $\theta_{\mathrm a}$ は、空気からコアへの屈折（$\sin\theta_{\mathrm a}=n_1\sin r$、$r$ はコアの中の軸からの角）と、境での臨界角（$\sin(\ang{90}-r)=\cos r=n_2/n_1$）から、
\begin{bkhenkei}[0.5]
\sin\theta_{\mathrm a}=n_1\sin r=n_1\sqrt{1-\cos^2r} & $\sin^2r+\cos^2r=1$（4-6） \\
=n_1\sqrt{1-\frac{n_2^{\,2}}{n_1^{\,2}}}=\sqrt{n_1^{\,2}-n_2^{\,2}} & 境で $\cos r=n_2/n_1$ \\
\end{bkhenkei}
です。この $\sqrt{n_1^{\,2}-n_2^{\,2}}$ を\textbf{開口数}（NA）といい、ファイバが光を受け入れる角の広さを表します。\bkidx{かいこうすう}{開口数（NA）}

\begin{phmodel}[光線で考える]
\phignore{光が波であること（回折・干渉。波-13）。境で一部がはね返ること（屈折するときも、少しは反射する）。物に吸収されて弱まること。}
\phassume{光は、同じ物の中ではまっすぐ進む線。境では、反射と屈折の法則にしたがう。}
\phlimit{光が通るすき間や物の大きさが、光の波長（約 \qtyrange{0.4}{0.8}{\micro m}）と同じくらいに小さいとき（波-13 の回折）。とても細いファイバ（コアが数 \si{\micro m}）では、光を波として扱う。}
\end{phmodel}

\Needspace{10\baselineskip}\bksubsection{5}{プリズムで光が曲がる}
\label{ph:波-12:sec5}
三角柱の形をしたガラスを\textbf{プリズム}といいます。\bkidx{ぷりずむ}{プリズム}光は、プリズムに入るときと出るときの2回屈折して、どちらも底の側へ曲がります。入る前と出たあとの光の向きの差を\textbf{偏角} $\delta$（へんかく）、プリズムの2つの面のなす角を\textbf{頂角} $\alpha$ といいます。

\noindent\begin{minipage}[c]{0.5\linewidth}\centering
\begin{tikzpicture}[x=1mm, y=1mm]
  \fill[bkBlue!12] (0,0) -- (40,0) -- (20,34.64) -- cycle;
  \draw[bkBlue!70!black, line width=0.8pt] (0,0) -- (40,0) -- (20,34.64) -- cycle;
  \rxmray{(-10.3,8.62)}{(8.66,15)}
  \draw[bkRed, line width=1pt] (8.66,15) -- (31.34,15);
  \draw[bkRed, line width=1pt, -{Stealth[length=2.2mm]}] (31.34,15) -- (50.3,8.62);
  \draw[bkGray, line width=0.5pt, dashed] (8.66,15) -- (38,24.9);
  \draw[bkGray, line width=0.5pt, dashed] (31.34,15) -- (20,18.82);
  \draw[bkInk, line width=0.5pt] ($(20,18.82)+(18.6:15)$) arc (18.6:-18.6:15);
  \node[font=\sffamily\footnotesize, anchor=west] at (35.6,18.8) {$\delta$};
  \draw[bkInk, line width=0.5pt] (17.5,29.9) arc (240:300:5);
  \node[font=\sffamily\footnotesize, anchor=north] at (20,27.6) {$\alpha$};
\end{tikzpicture}
\end{minipage}\hfill
\begin{minipage}[c]{0.46\linewidth}\raggedright
頂角\ang{60}、$n=1.5$ のプリズムの偏角を、屈折の法則を2回使って計算すると、右の表のようになります。入射角を変えると偏角も変わり、光がプリズムの中を\textbf{底に平行に、左右対称に通る}とき（入射角 約\ang{48.6}）に、偏角はいちばん小さく、約\ang{37.2}になります（\textbf{最小偏角}）。
\end{minipage}

\begin{bkcenter}
\begin{bktable}{colspec={Q[l,m,wd=24mm] *{7}{Q[c,m,wd=11mm]}}, rowsep=1.5pt, row{1}={bg=bkTheme!10, font=\sffamily\small}}
入射角〔°〕 & 30 & 40 & 45 & 48.6 & 55 & 60 & 70 \\
偏角 $\delta$〔°〕 & 47.1 & 38.5 & 37.4 & 37.2 & 37.7 & 38.9 & 42.9 \\
\end{bktable}
\bktabcaption{頂角\ang{60}・$n=1.5$ のプリズムの偏角（Python で計算）。入射角 約\ang{48.6} で最小。}\label{ph:波-12:tab-prism}
\end{bkcenter}

\noindent 最小偏角のときは、光が対称に通るので、中での屈折角は $\alpha/2$、入射角は $(\alpha+\delta)/2$ です。屈折の法則から、
\pheq{n=\frac{\sin\dfrac{\alpha+\delta_{\min}}{2}}{\sin\dfrac{\alpha}{2}}}{1}
となり、最小偏角をはかれば屈折率が求まります（ガラスの屈折率をはかる方法）。

\phboxhead{頂角の小さなプリズム}
頂角 $\alpha$ が小さいプリズムでは、入射角も小さいので、$\sin\theta\approx\theta$（5-17）とおけます。入る面で $\theta_{\text{入}}=nr_1$、出る面で $\theta_{\text{出}}=nr_2$、中の2つの屈折角の和は $r_1+r_2=\alpha$（三角形の角の関係）なので、
\[\delta=\theta_{\text{入}}+\theta_{\text{出}}-\alpha=n(r_1+r_2)-\alpha=(n-1)\,\alpha\]
です。\textbf{頂角の小さなプリズムは、入射角によらず、光を $(n-1)\alpha$ だけ曲げます}。頂角\ang{5}・$n=1.5$ なら約\ang{2.5}です。この式が、⑥でレンズを考える手がかりになります。

\Needspace{10\baselineskip}\bksubsection{6}{レンズと鏡の像}
\label{ph:波-12:sec6}
\label{ph:波-12:zou}まん中が厚いレンズを\textbf{凸レンズ}、まん中が薄いレンズを\textbf{凹レンズ}といいます。\bkidx{とつれんず}{凸レンズ}凸レンズに、軸（レンズの中心を通り、面に垂直な線）に平行な光を当てると、レンズを出た光は1点に集まります。この点を\textbf{焦点}、レンズの中心から焦点までの距離を\textbf{焦点距離} $f$\bkidx{しょうてんきょり}{焦点距離} といいます。\bkidx{しょうてん}{焦点}焦点は、レンズの両側に1つずつあります。

\phboxhead{レンズは、高さに比例して光を曲げる}
レンズを、細い輪切りの集まりと考えます。軸から高さ $y$ の所のレンズは、両面がわずかに傾いていて、頂角の小さなプリズムのように働きます。レンズの面が半径 $R_1$・$R_2$ の球の一部なら、高さ $y$ の所の2つの面のなす角は、軸に近いとき、およそ $y\,(1/R_1-1/R_2)$ です（面の傾きは「高さ÷半径」）。⑤の $(n-1)\alpha$ から、そこを通る光は
\[\delta=(n-1)\Bigl(\frac{1}{R_1}-\frac{1}{R_2}\Bigr)y\]
だけ、軸の側へ曲がります。曲がる角は、\textbf{高さ $y$ に比例}します。軸に平行に入った光は、高さ $y$ でも、$\delta$ 曲がって軸と交わるので、交わるまでの距離は $y/\delta$ で、$y$ によりません。だから、どの高さの光も1点（焦点）に集まるのです。その距離が焦点距離です。
\pheq{\frac{1}{f}=(n-1)\Bigl(\frac{1}{R_1}-\frac{1}{R_2}\Bigr)\qquad\delta=\frac{y}{f}}{1/m}
これを\textbf{レンズメーカーの式}といいます。\bkidx{れんずめーかーのしき}{レンズメーカーの式}$R$ は、面の中心がレンズより右（光が進む側）にあれば正、左にあれば負とします。両面が同じ半径 $R$ の凸レンズ（$R_1=R$、$R_2=-R$）では $1/f=2(n-1)/R$ で、$n=1.5$ なら $f=R$ です（例題7）。

\begin{phbaisuru}{レンズの屈折率と周りの屈折率の差}
レンズが曲げる力 $1/f$ は $(n-1)$ に比例します。$n-1$ は、本当は「レンズの屈折率÷周りの屈折率－1」です。空気の中で $1.5-1=0.5$ のレンズも、水の中では $1.5/1.333-1\approx0.125$ で、曲げる力が約4分の1、焦点距離が約4倍になります。水の中で目をあけると景色がぼやけるのは、目の表面（角膜）が、水の中では光をほとんど曲げなくなるからです。
\end{phbaisuru}

\phboxhead{レンズの式}
\label{ph:波-12:lens}軸の上の点（物）から、レンズまでの距離を $a$ とします。物から出て、レンズの高さ $y$ の所に来る光の傾きは $y/a$（小さな角）です。レンズで $y/f$ だけ軸の側へ曲がるので、出た光の傾きは $y/a-y/f$ です（下向きを負）。この光が軸と交わる点までの距離を $b$ とすると、
\begin{bkhenkei}[0.5]
y+b\Bigl(\frac{y}{a}-\frac{y}{f}\Bigr)=0 & $b$ 進むと高さが0になる \\
1+\frac{b}{a}-\frac{b}{f}=0 & 両辺を $y$ でわる（$y$ が消える） \\
\frac{1}{a}+\frac{1}{b}=\frac{1}{f} & 両辺を $b$ でわって整理 \\
\end{bkhenkei}
$y$ が消えたので、物の点から出た光は、どの高さを通っても、レンズの反対側の同じ点に集まります。そこに\textbf{像}ができます。\bkidx{ぞう}{像}
\pheq{\frac{1}{a}+\frac{1}{b}=\frac{1}{f}\qquad\text{倍率}=\frac{b}{a}}{1/m ＝ 1/m}
これが\textbf{レンズの式}です。\bkidx{れんずのしき}{レンズの式}レンズの中心を通る光は曲がらない（$y=0$）ので、物と像は中心をはさんで相似な三角形をつくり（4-3）、像の大きさは物の $b/a$ 倍（さかさま）です。

\Needspace{8\baselineskip}\begin{phflow}[レンズの式]
\phflowrow{$n_1\sin\theta_1=n_2\sin\theta_2$（③）}{薄いレンズ・軸に近い光\\$\sin\theta\approx\theta$}{高さ $y$ の光を $\dfrac yf$ 曲げる}
\phflowrow{高さ $y$ の光を $\dfrac yf$ 曲げる}{物から $a$、像まで $b$}{$\dfrac1a+\dfrac1b=\dfrac1f$、倍率 $\dfrac ba$}
\end{phflow}

\begin{bkfigure}
\begin{tikzpicture}[x=0.85mm, y=1mm]
  \draw[bkGray, line width=0.5pt, dashed] (-50,0) -- (32,0);
  \rxmlens{0}{16}
  \fill (15,0) circle (0.8) (-15,0) circle (0.8);
  \node[font=\sffamily\footnotesize, anchor=south west] at (16,0.6) {F$'$};
  \node[font=\sffamily\footnotesize, anchor=north] at (-15,-0.8) {F};
  \rxmobj{(-45,0)}{12}{bkGreen!60!black}
  \draw[bkOrange!80!black, line width=1.6pt, -{Stealth[length=2.6mm, width=2.2mm]}] (22.5,0) -- (22.5,-6);
  \draw[bkRed, line width=0.9pt] (-45,12) -- (0,12) -- (27,-9.6);
  \draw[bkBlue!70!black, line width=0.9pt] (-45,12) -- (27,-7.2);
  \draw[bkPurple, line width=0.9pt] (-45,12) -- (0,-6) -- (27,-6);
  \node[font=\sffamily\footnotesize, anchor=east] at (-46,6) {物};
  \node[font=\sffamily\footnotesize, anchor=north] at (22.5,-10.5) {像（実像）};
  \node[anchor=west, font=\small, text width=42mm, align=left] at (40,6) {①軸に平行な光→F$'$ を通る\par ②中心を通る光→まっすぐ\par ③F を通る光→軸に平行\par 3本は1点に集まる（$a=3f$ のとき $b=1.5f$、倍率0.5）};
\end{tikzpicture}
\bkfigcaption{凸レンズの像の作図。物の先から出る3本の光（①上の実線、②まん中、③下）を引けば、像の位置と大きさがわかる。}\label{ph:波-12:fig-sakuzu}
\end{bkfigure}

\phboxhead{実像と虚像}
物が焦点より遠い（$a>f$）とき、$b$ は正で、レンズの反対側に、光が本当に集まった像ができます（\textbf{実像}。さかさま）。\bkidx{じつぞう}{実像}スクリーンに映せます。物が焦点より近い（$a<f$）とき、レンズの式の $b$ は負になります。光は集まらずに広がって出ていきますが、目には、レンズの手前（物と同じ側）の $|b|$ の所から光が出ているように見えます。これが\textbf{虚像}で、正立で大きく見えます（虫めがね）。\bkidx{きょぞう}{虚像}

\begin{bkfigure}
\begin{tikzpicture}[x=1mm, y=1mm]
  \draw[bkGray, line width=0.5pt, dashed] (-36,0) -- (40,0);
  \rxmlens{0}{16}
  \fill (15,0) circle (0.8) (-15,0) circle (0.8);
  \node[font=\sffamily\footnotesize, anchor=north] at (15,-0.8) {F$'$};
  \node[font=\sffamily\footnotesize, anchor=north] at (-15,-0.8) {F};
  \rxmobj{(-10,0)}{5}{bkGreen!60!black}
  \draw[bkOrange!60, line width=1.6pt, dashed, -{Stealth[length=2.6mm, width=2.2mm]}] (-30,0) -- (-30,15);
  \draw[bkRed, line width=0.9pt] (-10,5) -- (0,5) -- (30,-5);
  \draw[bkBlue!70!black, line width=0.9pt] (-10,5) -- (20,-10);
  \draw[bkRed!50, line width=0.6pt, dashed] (0,5) -- (-30,15);
  \draw[bkBlue!50, line width=0.6pt, dashed] (-10,5) -- (-30,15);
  \rxmeye{(38,-4)}{180}
  \node[font=\sffamily\footnotesize, anchor=north] at (-10,-0.8) {物};
  \node[font=\sffamily\footnotesize, anchor=east] at (-31,9) {虚像};
  \node[anchor=west, font=\small, text width=50mm, align=left] at (46,6) {物が焦点の内側（$a=10$、$f=15$）。\par 出た光を逆にのばす（破線）と、$|b|=30$ の所で交わる。\par 正立で3倍の虚像。};
\end{tikzpicture}
\bkfigcaption{虫めがねの虚像。レンズを出た光（実線）は広がるが、目は破線の先の1点から来たと感じる。}\label{ph:波-12:fig-kyozou}
\end{bkfigure}

\begin{phbaisuru}{物までの距離 $a$}
$a$ を $f$ に近い所から大きくしていくと、像の距離 $b$ と倍率 $b/a$ は次のように変わります（$f=\qty{10}{cm}$）。
\begin{bkcenter}
\begin{bktable}{colspec={Q[l,m,wd=20mm] *{6}{Q[c,m,wd=11mm]}}, rowsep=1.5pt, row{1}={bg=bkTheme!10, font=\sffamily\small}}
$a$〔cm〕 & 5 & 15 & 20 & 30 & 60 & 200 \\
$b$〔cm〕 & $-10$ & 30 & 20 & 15 & 12 & 10.5 \\
倍率 $b/a$ & 2（正立） & 2 & 1 & 0.5 & 0.2 & 0.053 \\
\end{bktable}
\end{bkcenter}
$a$ を2倍（\qty{30}{cm}→\qty{60}{cm}）にすると、$b$ は\qty{15}{cm}から\qty{12}{cm}へ少し縮み、像は0.5倍から0.2倍へ小さくなります。遠い物ほど、像は焦点の近く（$b\approx f$）に小さくできます。$a=2f$ のとき、同じ大きさのさかさまの像が $b=2f$ にできます。
\end{phbaisuru}

\phboxhead{凹レンズ}
凹レンズは、軸に平行な光を、手前の焦点から広がって出てくるように曲げます。レンズの式で $f$ を負にすれば、そのまま使えます。$f=\qty{-20}{cm}$ のレンズの\qty{30}{cm}前に物を置くと、$1/b=1/(-20)-1/30$ より $b=\qty{-12}{cm}$、手前\qty{12}{cm}に、0.4倍の正立の虚像が見えます。近視のめがね（⑨）に使います。

\phboxhead{球面鏡}
\label{ph:波-12:kyumen}内側が鏡の、球の一部の形をした鏡を\textbf{凹面鏡}といいます。\bkidx{おうめんきょう}{凹面鏡}半径 $R$ の凹面鏡では、高さ $y$ の所の鏡の面は、およそ $y/R$ 傾いています。反射の法則（②）で、光は面の傾きの2倍、$2y/R$ 曲げられます。レンズの $y/f$ と比べると、凹面鏡は焦点距離 $f=R/2$ のレンズと同じ働きをします（光は反対向きにもどる）。
\pheq{\frac{1}{a}+\frac{1}{b}=\frac{2}{R}=\frac{1}{f}\qquad f=\frac R2}{1/m}
これを\textbf{球面鏡の式}といいます（像の距離 $b$ は、鏡の前を正とする）。外側が鏡の\textbf{凸面鏡}は $f=-R/2$ です。道路のカーブミラー（凸面鏡、$R=\qty{1}{m}$）で\qty{2}{m}先の車を見ると、$1/b=-2-0.5$ より $b=\qty{-0.4}{m}$、倍率0.2の小さな虚像で、広い範囲が映ります。

\Needspace{10\baselineskip}\bksubsection{7}{行列で光線を追う}
\label{ph:波-12:sec7}
カメラのレンズは、何枚ものレンズを組み合わせています。1枚ずつ作図していくのは大変です。7-3 では、光線を（高さ $y$、傾き $\theta$）の組で表し、すき間とレンズを行列で表しました。
\[\text{すき間 }d：\begin{pmatrix}1&d\\0&1\end{pmatrix}\qquad\text{レンズ }f：\begin{pmatrix}1&0\\-1/f&1\end{pmatrix}\]
すき間では、高さが $d\theta$ 変わり（傾きは同じ）、レンズでは、傾きが $y/f$ 減ります（高さは同じ）。⑥で作った「レンズは高さに比例して曲げる」が、そのまま行列になっています。光が順に通る要素の行列を、\textbf{右から左へ}かけていけば（7-3 の「（レンズ）×（すき間）」）、全体が1つの行列になります。

\phboxhead{行列から、レンズの式を導き直す}
物からレンズまで $a$、レンズから $b$ 進むと、全体の行列は
\begin{bkhenkei}[0.5]
\begin{pmatrix}1&b\\0&1\end{pmatrix}\begin{pmatrix}1&0\\-1/f&1\end{pmatrix}\begin{pmatrix}1&a\\0&1\end{pmatrix} & すき間 $a$→レンズ→すき間 $b$ \\
=\begin{pmatrix}1-b/f&a+b-ab/f\\-1/f&1-a/f\end{pmatrix} & 7-3 の積（行と列の内積） \\
\end{bkhenkei}
です。この行列の成分を、左上から $A$、$B$、$C$、$D$ とよびます。出口の高さは $y_{\text{出}}=Ay_{\text{入}}+B\theta_{\text{入}}$ です。物の1点から出た光が、傾き $\theta_{\text{入}}$ によらず同じ点に集まる（像ができる）条件は、$B=0$ です。
\[a+b-\frac{ab}{f}=0\quad\text{両辺を}\ ab\ \text{でわって}\quad\frac1b+\frac1a=\frac1f\]
⑥のレンズの式が出ました。そのとき $y_{\text{出}}=Ay_{\text{入}}$ なので、$A=1-b/f=-b/a$ が倍率です（負の符号はさかさま）。

\phboxhead{2枚のレンズ}
焦点距離 $f_1$ のレンズの $d$ 後ろに、$f_2$ のレンズを置きます。全体の行列は（2枚め）×（すき間）×（1枚め）で、その左下の成分 $C$ から、2枚を1枚のレンズと見たときの焦点距離（\textbf{合成焦点距離}）がわかります。
\pheq{\frac1f=\frac1{f_1}+\frac1{f_2}-\frac{d}{f_1f_2}}{1/m}
$d=0$（くっつける）なら、$1/f$ は和になります。めがねの「度数」（⑨）がたし算できるのは、このためです。

\begin{phcode}[Python]
import numpy as np
def gap(d):  return np.array([[1, d], [0, 1]])      \# すき間 d〔m〕
def lens(f): return np.array([[1, 0], [-1/f, 1]])   \# レンズ f〔m〕
M = lens(0.1) @ gap(0.05) @ lens(0.1)   \# 右から順に光が通る
print(np.round(M, 3))
print("合成の焦点距離", round(-1/M[1, 0], 4))
ray = np.array([0.01, 0.0])             \# 高さ 0.01 m、軸に平行
y, th = M @ ray
print("出た光", round(y, 4), round(th, 4), "軸まで", round(-y/th, 4))
\end{phcode}
\begin{phcodeout}[実行結果]
[[  0.5    0.05]
 [-15.     0.5 ]]
合成の焦点距離 0.0667
出た光 0.005 -0.15 軸まで 0.0333
\end{phcodeout}
\noindent \texttt{@} は numpy の行列の積です。$f=\qty{0.1}{m}$ のレンズ2枚を\qty{0.05}{m}はなすと、合成焦点距離は約\qty{0.067}{m}で、平行な光は2枚めのレンズの\qty{0.033}{m}後ろに集まります（例題9）。レンズの枚数がふえても、行列をかけるだけです。実際のカメラの設計では、何万本もの光線を、この考えで（行列ではなく、正確な屈折の法則で）追いかけます。

\Needspace{10\baselineskip}\bksubsection{8}{分散：色で曲がり方が違う、虹}
\label{ph:波-12:sec8}
\label{ph:波-12:bunsan}1666年ごろ、ニュートンは、太陽の光をプリズムに通すと、赤・だいだい・黄・緑・青・紫の色の帯に分かれることを確かめました。分かれた1色の光を、もう一度プリズムに通しても、それ以上は分かれません。白い光は、いろいろな色の光がまざったものなのです（熱-13 でも、光が波長ごとに分けられることを使いました）。

色によって分かれるのは、\textbf{屈折率が光の色（波長）で少しずつ違う}からです。ガラスや水では、紫の光（波長が短い）ほど屈折率が大きく、赤の光（波長が長い）ほど小さくなります。だから、\textbf{プリズムでは、紫の光ほどよく曲がり、赤の光は曲がりにくい}のです。このように、屈折率が波長で変わることを\textbf{分散}（ぶんさん）といい、色の帯を\textbf{スペクトル}といいます。\bkidx{ぶんさん}{分散（光）}\bkidx{すぺくとる}{スペクトル（光）}

\begin{bkfigure}
\begin{tikzpicture}[x=1mm, y=1mm]
  \fill[bkBlue!10] (0,0) -- (40,0) -- (20,34.64) -- cycle;
  \draw[bkBlue!70!black, line width=0.8pt] (0,0) -- (40,0) -- (20,34.64) -- cycle;
  \draw[bkInk!60, line width=2.2pt] (-14,7.9) -- (8.66,15);
  \foreach \yy/\aa/\cc in {16/15/bkRed, 15.5/16.6/bkOrange, 15/18.2/bkYellow!80!black, 14.6/19.8/bkGreen!70!black, 14.1/21.4/bkBlue, 13.6/23/bkPurple} {
    \pgfmathsetmacro\xx{40-0.5774*\yy}
    \draw[\cc, line width=0.8pt] (8.66,15) -- (\xx,\yy);
    \draw[\cc, line width=1.1pt] (\xx,\yy) -- ++(-\aa:34);}
  \node[font=\sffamily\footnotesize, anchor=west, text=bkRed!80!black] at (64.5,7.2) {赤（曲がりにくい）};
  \node[font=\sffamily\footnotesize, anchor=west, text=bkPurple] at (61,-1.6) {紫（よく曲がる）};
  \node[font=\sffamily\footnotesize, anchor=south] at (-14,8.4) {白い光};
\end{tikzpicture}
\bkfigcaption{プリズムによる分散（曲がり方の差は大きくかいた）。屈折率は紫ほど大きいので、紫がいちばん大きく曲がり、色の帯（スペクトル）に分かれる。}\label{ph:波-12:fig-bunsan}
\end{bkfigure}

\begin{bkcenter}
\begin{bktable}{colspec={Q[l,m,wd=36mm] Q[c,m,wd=34mm] Q[c,m,wd=34mm]}, rowsep=1.5pt, row{1}={bg=bkTheme!10, font=\sffamily\small}}
 & 赤（\qty{656}{nm}） & 紫（\qty{405}{nm}付近） \\
水（\qty{20}{\celsius}）の屈折率 & 1.331 & 1.343 \\
\end{bktable}
\bktabcaption{水の屈折率の、色による違い（可視光は約 \qtyrange{380}{780}{nm}、熱-13・波-13）。差は1\,\%ほどだが、虹の色を分けるには十分。}\label{ph:波-12:tab-mizu}
\end{bkcenter}

\phboxhead{虹の角}
雨上がりに、太陽を背にして立つと、虹が見えます。空に浮かぶ小さな水の粒（ほぼ球）の中で、光が屈折・反射・屈折して、もどってくるのです（図\ref{ph:波-12:fig-niji}）。

\begin{bkfigure}
\begin{tikzpicture}[x=1mm, y=1mm]
  \fill[bkBlue!10] (0,0) circle (20);
  \draw[bkBlue!70!black, line width=0.8pt] (0,0) circle (20);
  \draw[bkGray, line width=0.4pt, dashed] (0,0) -- (-10.1,17.2) (0,0) -- (18.6,7.3) (0,0) -- (4.2,-19.6);
  \rxmray{(-46,17.2)}{(-10.1,17.2)}
  \draw[bkRed, line width=1pt] (-10.1,17.2) -- (18.6,7.3) -- (4.2,-19.6);
  \draw[bkRed, line width=1pt, -{Stealth[length=2.2mm]}] (4.2,-19.6) -- (-17.9,-39.8);
  \draw[bkGray, line width=0.5pt, dashed] (4.2,-19.6) -- (-26,-19.6);
  \draw[bkInk, line width=0.5pt] (-10,-19.6) arc (180:222.4:14.2);
  \node[font=\sffamily\footnotesize, anchor=east] at (-11,-24) {約\ang{42}};
  \node[font=\sffamily\footnotesize, anchor=south] at (-36,17.7) {太陽の光};
  \node[font=\sffamily\footnotesize, anchor=west] at (19,8) {反射};
  \node[anchor=west, font=\small, text width=48mm, align=left] at (28,-10) {入る：屈折（$\theta_1\to\theta_2$）\par 奥の面：反射\par 出る：屈折\par もどる光は、太陽の光と反対の向きから約\ang{42}の向きに集まる。};
\end{tikzpicture}
\bkfigcaption{水の粒の中の光（赤い光、粒の中心からの高さで決まる入射角 $\theta_1$ が約\ang{59.5}のとき）。破線は粒の中心と光の当たる点を結ぶ線（法線）。}\label{ph:波-12:fig-niji}
\end{bkfigure}

\noindent 入射角を $\theta_1$、粒の中での屈折角を $\theta_2$（$\sin\theta_1=n\sin\theta_2$）とすると、入るとき・反射のとき・出るときの3回の向きの変化から、もどってくる光が、太陽の光と正反対の向き（自分の頭の影の向き）となす角は
\pheq{\phi=4\theta_2-2\theta_1}{°}
です。いろいろな高さに当たる光で計算すると（$n=1.333$）、$\theta_1=\ang{40}$ で\ang{35.3}、\ang{55} で\ang{41.7}、\ang{60} で\ang{42.1}、\ang{70} で\ang{39.3}と、$\phi$ は約\ang{42}より大きくなりません。いちばん大きい所の近くでは、$\theta_1$ が少し変わっても $\phi$ がほとんど変わらないので、たくさんの光がこの角に集まり、明るく見えます。$\dd\phi/\dd\theta_1=0$（5-7）から、そのときの入射角は $\cos^2\theta_1=(n^2-1)/3$ を満たします（例題10）。

\begin{bkfigure}
\begin{tikzpicture}
\begin{axis}[bkplot, width=104mm, height=50mm, xmin=0, xmax=95, ymin=0, ymax=48, xtick={20,40,60,80,90}, ytick={10,20,30,40}, yticklabel style={/pgf/number format/fixed}, xlabel={入射角 $\theta_1$（°）}, ylabel={もどる角 $\phi$（°）}]
  \addplot[bkRed, line width=1.4pt, domain=0:90, samples=120] {4*asin(sin(x)/1.331)-2*x};
  \addplot[bkPurple, line width=1.4pt, dash pattern=on 4pt off 2pt, domain=0:90, samples=120] {4*asin(sin(x)/1.343)-2*x};
  \node[font=\sffamily\footnotesize, anchor=south, text=bkRed!80!black] at (axis cs:72,43.5) {赤（実線）\ang{42.4}};
  \node[font=\sffamily\footnotesize, anchor=north, text=bkPurple] at (axis cs:56,22) {紫（破線）\ang{40.6}};
\end{axis}
\end{tikzpicture}
\bkfigcaption{入射角ともどる角。どちらの色も、いちばん大きい角（赤 \ang{42.4}、紫 \ang{40.6}）の近くに光が集まる。}\label{ph:波-12:fig-nijikaku}
\end{bkfigure}

\noindent 赤（$n=1.331$）は約\ang{42.4}、紫（$n=1.343$）は約\ang{40.6}に集まります。頭の影の向きから\ang{42.4}の円の上の粒からは赤い光が、\ang{40.6}の円の上の粒からは紫の光が目に届くので、\textbf{虹は外側が赤、内側が紫}の円になります。粒の中で2回反射した光が作る、外側の薄い虹（副虹、約\ang{50}〜\ang{53}）は、色の並びが反対です。

\Needspace{10\baselineskip}\bksubsection{9}{光学機器：目・カメラ・虫めがね・望遠鏡}
\label{ph:波-12:sec9}

\phboxhead{目とめがね}
目は、角膜と水晶体（レンズ）で、網膜の上に実像を結ぶ光学機器です。遠くを見るときと近くを見るときで、物までの距離 $a$ が変わっても、網膜までの距離 $b$ は変えられないので、水晶体の厚さを変えて焦点距離 $f$ を調節します（レンズの式）。近くの物がはっきり見えるいちばん近い距離は、年れいとともに遠くなります（老眼）。

レンズの曲げる力を、焦点距離〔m〕の逆数で表したものを\textbf{度数}（ディオプトリー、D）といいます。$f=\qty{0.5}{m}$ の凸レンズは $+2$\,D、$f=\qty{-0.5}{m}$ の凹レンズは $-2$\,D です。\bkidx{どすう}{度数（レンズ）}遠くの物の像が網膜より前にできる近視では、凹レンズで光を少し広げてから目に入れます。くっつけたレンズの度数はたし算できる（⑦）ので、めがねの度数は「目に足りない分」を表します。

\phboxhead{カメラ}
カメラは、レンズで受光素子（スマホではイメージセンサー）の上に、さかさまの実像を作ります。物までの距離が変わると、レンズの式の $b$ が変わるので、レンズを前後に動かしてピントを合わせます（例題8）。

\phboxhead{虫めがね}
物を目に近づけると大きく見えますが、約\qty{25}{cm}（\textbf{明視の距離}）より近いと、ピントが合いません。虫めがね（焦点距離 $f$）を目の前に置き、物を焦点の近くに置くと、物から出た光はほぼ平行になって目に入り、目は楽に像を見られます。そのとき、物が目に見える角は、明視の距離で見るときの約 $25/f$ 倍（$f$ は cm）に大きくなります。
\pheq{\text{虫めがねの倍率}\approx\frac{\qty{25}{cm}}{f}}{1}
（像を明視の距離にできるようにすると、$25/f+1$ 倍。例題11）。

\phboxhead{顕微鏡と望遠鏡}
\textbf{顕微鏡}は、焦点距離の短い対物レンズで、小さな物の大きな実像を作り、それを接眼レンズ（虫めがね）で見ます。全体の倍率は、2つの倍率の積です（たとえば40倍×10倍＝400倍）。\textbf{望遠鏡}は、焦点距離の長い対物レンズで、遠くの物の実像を焦点の所に作り、それを焦点距離の短い接眼レンズで見ます。物が見える角は、
\pheq{\text{望遠鏡の倍率}=\frac{f_{\text{対物}}}{f_{\text{接眼}}}}{1}
倍になり、筒の長さはおよそ $f_{\text{対物}}+f_{\text{接眼}}$ です（例題12）。

\begin{bkfigure}
\begin{tikzpicture}[x=1mm, y=1mm]
  \draw[bkGray, line width=0.5pt, dashed] (-14,0) -- (98,0);
  \rxmlens{0}{12}
  \rxmlens{90}{6}
  \fill (80,0) circle (0.8);
  \draw[bkRed, line width=0.9pt] (-14,10) -- (0,10) -- (80,0) -- (90,-1.25) -- (104,-1.25);
  \draw[bkRed, line width=0.9pt] (-14,-10) -- (0,-10) -- (80,0) -- (90,1.25) -- (104,1.25);
  \draw[bkInk, line width=0.5pt, {Stealth[length=2pt]}-{Stealth[length=2pt]}] (0,-15) -- (80,-15);
  \draw[bkInk, line width=0.5pt, {Stealth[length=2pt]}-{Stealth[length=2pt]}] (80,-15) -- (90,-15);
  \node[font=\sffamily\footnotesize, anchor=north] at (40,-15.5) {$f_{\text{対物}}$};
  \node[font=\sffamily\footnotesize, anchor=north] at (85,-15.5) {$f_{\text{接眼}}$};
  \node[font=\sffamily\footnotesize, anchor=south] at (0,12.5) {対物レンズ};
  \node[font=\sffamily\footnotesize, anchor=south] at (90,12.5) {接眼レンズ};
\end{tikzpicture}
\bkfigcaption{望遠鏡（ケプラー式）。軸に平行に入った光（実線。遠くの1点からの光）は、共通の焦点を通ってから接眼レンズで再び平行になり、目に入る。ここでは軸にそった光だけをかいた。}\label{ph:波-12:fig-bouenkyou}
\end{bkfigure}

\Needspace{10\baselineskip}\bksubsection{10}{明るさをはかる：光束・光度・照度}
\label{ph:波-12:sec10}
部屋の照明や太陽の光の「明るさ」を数で表すには、いくつかの量を使い分けます。どれも、人の目が感じる明るさに合わせて決められた量です。

\begin{bkcenter}
\begin{bktable}{colspec={Q[l,m,wd=18mm] Q[c,m,wd=12mm] Q[l,m,wd=50mm] Q[l,m,wd=36mm]}, rowsep=1.5pt, row{1}={bg=bkTheme!10, font=\sffamily\small}}
量 & 単位 & 意味 & 例 \\
光束 $\Phi$ & lm & 光源が1秒に出す光の量（目の感じ方で重みをつけた） & LED 電球 60形相当 約810 lm \\
光度 $I$ & cd & ある向きへ、立体角 1 sr あたりに出す光束 & ろうそく1本 約1 cd \\
照度 $E$ & lx & 照らされる面の 1 m$^2$ に届く光束 & 事務室 750 lx \\
\end{bktable}
\bktabcaption{測光の量（ルーメン lm、カンデラ cd、ルクス lx。sr は立体角の単位ステラジアン、8-8）。}\label{ph:波-12:tab-sokkou}
\end{bkcenter}

\phboxhead{照度は距離の2乗に反比例する}
どの向きにも光度 $I$（波-7・13 の強さ $I$〔W/m$^2$〕とは別の量）で光を出す小さな光源は、全部で $\Phi=4\pi I$ の光束を出します（球の全体の立体角は $4\pi$、8-8）。この光束は、半径 $r$ の球の面（面積 $4\pi r^2$）に一様に広がるので、距離 $r$ の所で光に垂直な面の照度は
\pheq{E=\frac{4\pi I}{4\pi r^2}=\frac{I}{r^2}\qquad\text{面が}\ \theta\ \text{傾くと}\quad E=\frac{I\cos\theta}{r^2}}{lx ＝ $\dfrac{\text{cd}}{\text{m}^2}$}
\label{ph:波-12:gyaku}です（$\theta$ は、面の法線と光のなす角。面が傾くと、同じ光束が広い面積に広がる）。これを\textbf{照度の逆2乗の法則}といいます。\bkidx{しょうどのぎゃく2じょうのほうそく}{照度の逆2乗の法則}

\begin{phbaisuru}{光源までの距離 $r$}
照度は $\tfrac14$ になります。光束が4倍の面積に広がるからです。光度 \qty{1000}{cd}の照明から\qty{2}{m}で\qty{250}{lx}、\qty{4}{m}で約\qty{63}{lx}です。本を読むときに明かりに近づくと、ずっと明るくなるのはこのためです。
\end{phbaisuru}

\begin{phmijika}[部屋の明るさの目安]
日本の照明の規格（JIS Z 9110）では、事務所の事務室で\qty{750}{lx}を目安としています。晴れた日に太陽へ向けた面は約10万 lx で、部屋の明かりの100倍以上です。それでも目が慣れて同じくらいに感じるのは、目が明るさの「比」を感じるからです（音の大きさの dB と同じ考え。3-12・波-9）。
\end{phmijika}

\begin{phmada}[虫や鳥は、光のどんな性質で道を知る？]
ミツバチは、太陽が雲にかくれていても、青空の一部が見えれば、巣と花の間の向きがわかります。空の光には、ゆれる向きのかたより（偏光。波-13 ⑥）の模様があり、ミツバチはそれを目で感じて、太陽の向きを知るのです（フォン・フリッシュ、1949年）。その後、目のどの部分が偏光を感じるか、脳のどこ（中心複合体とよばれる所）に方角を表す細胞があるかまでは、わかってきました。でも、脳がその情報から、どのように進む向きを決めて体を動かしているのかは、まだはっきりしていません。渡り鳥が、太陽・星・地球の磁場などを、どのように組み合わせて道を知るのかも、研究が続いています。
\end{phmada}

\begin{phsaki}{光学設計と、波としての光}
この章では、光を線として、軸に近い光だけを扱いました（$\sin\theta\approx\theta$）。軸から遠い光や、ななめに入る光では、像がぼやけたり、ゆがんだり、色がにじんだりします（\textbf{収差}）。カメラのレンズが何枚もの組み合わせなのは、収差を打ち消すためです。光学設計の専門の本では、たくさんの光線をコンピュータで追いかけて、収差を小さくする方法を学びます。光が波であることから起きる回折や干渉は、波-13 で、光が電気と磁気の波であることは物理シリーズの『電磁気学』（電-15）で、光が粒（光子）でもあることは『原子』（原-2）で学びます。
\end{phsaki}

\begin{senmonsho}[光学の専門書の書き方]
\begin{bkhonbun}
近軸光線は $(h,\ u)$ で表され，転送行列は $\begin{pmatrix}1&d\\0&1\end{pmatrix}$，屈折行列は $\begin{pmatrix}1&0\\-\phi&1\end{pmatrix}$ である．ここで $\phi=1/f$ は屈折力である．系の $ABCD$ 行列において，結像条件は $B=0$，横倍率は $A$ である．
\end{bkhonbun}
\bkyaku{$h$ はこの本の高さ $y$、$u$ は傾き $\theta$。「近軸」は軸に近い光（$\sin\theta\approx\theta$）。$\phi$（ファイ）は屈折力で、度数（D）と同じもの。本によっては、光線を（傾き、高さ）の順に書いたり、屈折率をかけた傾き $nu$ を使ったりするので、行列の形が少し違う。}
\end{senmonsho}

\begin{taisetsu}
\def\tsl#1{\makebox[7.6\zw][l]{\sffamily\bfseries\small\color{\bkTT}#1}}
\tsl{反射の法則}入射角＝反射角。平面鏡の像は鏡に対して対称\\[3pt]
\tsl{屈折の法則}$n_1\sin\theta_1=n_2\sin\theta_2$、$n=c/v$（フェルマーの原理）\\[3pt]
\tsl{全反射}$\sin\theta_{\mathrm c}=n_2/n_1$。ファイバの NA $=\sqrt{n_1^{\,2}-n_2^{\,2}}$\\[3pt]
\tsl{レンズ}$\dfrac1a+\dfrac1b=\dfrac1f$、倍率 $\dfrac ba$。$\dfrac1f=(n-1)\Bigl(\dfrac1{R_1}-\dfrac1{R_2}\Bigr)$\\[3pt]
\tsl{行列}すき間・レンズの積。像は $B=0$、倍率は $A$\\[3pt]
\tsl{虹・機器}虹は赤 \ang{42}・紫 \ang{41}。望遠鏡 $f_{\text{対物}}/f_{\text{接眼}}$\\[3pt]
\tsl{照度}$E=I\cos\theta/r^2$（距離2倍で $\tfrac14$）
\end{taisetsu}

\begin{kotoba}
\small
\bkword{光線}{光の進む道を、まっすぐな線として表したもの。}
\bkword{反射の法則}{入射角と反射角が等しいこと。}
\bkword{屈折率}{真空の光速 $c$ が、物の中の光の速さ $v$ の何倍か。$n=c/v$。}
\bkword{フェルマーの原理}{光は、かかる時間がいちばん短い道（せいかくには、時間が極値になる道）を通る。}
\bkword{全反射・臨界角}{屈折率の小さい物へ出られずに、全部はね返ること。その境の入射角。}
\bkword{焦点距離}{レンズの中心から、平行な光が集まる点（焦点）までの距離。}
\bkword{実像・虚像}{光が本当に集まってできる像・光が集まって見えるだけの像。}
\bkword{分散}{屈折率が光の色（波長）で変わること。紫ほどよく曲がる。}
\bkword{照度}{面の 1 m$^2$ に届く光束。単位 lx。}
\end{kotoba}

\bkkigocards{\bkkigocard{$n$}{エヌ}{屈折率 $c/v$。単位なし。第4部で使う（波-6 ではモードの番号）}{refractive index}{波-12}}
{\bkkigocard{$E$}{イー}{照度〔lx〕。波-12 で使う（ヤング率・エネルギーは出ない）}{illuminance}{波-12}}

%% ======================== 例題 ========================
\Needspace{10\baselineskip}\bksection{reidai}{例題でたしかめよう}
\noindent 例題では、空気の屈折率を1、水を1.333、ガラスを1.5 とします。角は、電卓を DEG（度）にして計算し、3けたで書きます。

\label{ph:波-12:ex1}
\begin{phreidai}[全身を映す鏡]
身長\qty{160}{cm}の人の目は、頭のてっぺんより\qty{10}{cm}下にあります。かべにかけた鏡に、頭からつま先まで全身を映すには、鏡のたての長さは少なくとも何 cm いりますか。鏡の下のはしは、床から何 cm の高さにすればよいですか。
\phbamen
平面鏡で、自分の全身を見る場面です。
\phmushi
鏡のふちの厚さ、目の大きさ（目を1点とする）を無視します。人も鏡もまっすぐ立っているとします。
\phzu
床から目の高さ\qty{150}{cm}、頭\qty{160}{cm}、つま先\qty{0}{cm}。鏡までの距離を $L$ とする。
\phriyu
頭やつま先から出て、鏡で反射して目に届く光は、反射の法則（入射角＝反射角、②）にしたがいます。目と物の高さが同じ距離 $L$ で鏡をはさむので、反射する点は、2つの高さのちょうどまん中です。
\phtateru
\[\text{上のはし}=\frac{160+150}{2},\qquad\text{下のはし}=\frac{0+150}{2}\]
\phtoku
\[\text{上のはし}=\qty{155}{cm},\qquad\text{下のはし}=\qty{75}{cm},\qquad\text{長さ}=155-75=\qty{80}{cm}\]
\phtashikame
\qty{80}{cm}は身長\qty{160}{cm}のちょうど半分です。答えに $L$ が出てこないので、鏡から遠ざかっても、近づいても、必要な長さは変わりません。
\kotae 長さ\qty{80}{cm}、下のはしは床から\qty{75}{cm}
\end{phreidai}
\begin{jibun}
身長\qty{120}{cm}、目が頭のてっぺんより\qty{8}{cm}下にある子どもが、全身を映すのに要る鏡の長さと、鏡の下のはしの高さを求めましょう。
\bkjans{長さ $120\div2=\qty{60}{cm}$。下のはしは $112\div2=\qty{56}{cm}$（上のはしは $(120+112)\div2=\qty{116}{cm}$）}
\end{jibun}

\label{ph:波-12:ex2}
\begin{phreidai}[空気から水へ入る光]
空気から水面へ、入射角\ang{45}で光が入ります。屈折角を求めましょう。水の中の光の速さも求めましょう。
\phbamen
プールの水面に、ななめから光が差しこむ場面です。
\phmushi
水面で反射する光、水面の波立ちを無視します。
\phzu
③の図\ref{ph:波-12:fig-kussetsu}。$n_1=1$、$\theta_1=\ang{45}$、$n_2=1.333$。
\phriyu
2つの物の境での向きの変化は、屈折の法則（③）。光の速さは $v=c/n$（③）。
\phtateru
\[\sin\theta_2=\frac{n_1\sin\theta_1}{n_2},\qquad v=\frac{c}{n_2}\]
\phtoku
\[\sin\theta_2=\frac{1\times0.7071}{1.333}=0.5305,\qquad\theta_2=\sin^{-1}0.5305\approx\ang{32.0}\]
\[v=\frac{3.00\times10^8}{1.333}\approx\qty{2.25e8}{m/s}\]
\phtashikame
$\theta_2<\theta_1$ で、遅い水へ入ると法線に近づきます。図\ref{ph:波-12:fig-kakudo}の実線で、入射角\ang{45}の所を読むと、約\ang{32}です。
\kotae 約\ang{32}、約\qty{2.25e8}{m/s}
\end{phreidai}
\begin{jibun}
（1）空気からガラスへ、入射角\ang{30}で入る光の屈折角を求めましょう。（2）水の中から水面へ、入射角\ang{30}で当たった光が空気へ出るときの屈折角を求めましょう。
\bkjans{（1）$\sin\theta_2=0.5\div1.5=0.333$、約\ang{19.5}。（2）$\sin\theta_2=1.333\times0.5=0.667$、約\ang{41.8}}
\end{jibun}

\label{ph:波-12:ex3}
\begin{phreidai}[水の中のストローの先は、どこに見える？（はてな？）]
深さ\qty{10}{cm}の水の底に、ストローの先がとどいています。真上に近い向きから見ると、ストローの先は水面から何 cm の深さに見えますか。
\phbamen
はてな？の、折れて見えるストローの場面です。
\phmushi
コップのガラスの壁での屈折を無視し、水面の上から見るとします。目は、真上に近い向き（光の角が小さい）にあるとします。
\phzu
③の見かけの深さの図。$h=\qty{10}{cm}$、$n=1.333$、見かけの深さ $h'$。
\phriyu
小さな角では、屈折の法則から見かけの深さは $h'=h/n$（③の流れ図）です。
\phtateru
\[h'=\frac{h}{n}\]
\phtoku
\[h'=\frac{10}{1.333}\approx\qty{7.5}{cm}\]
\phtashikame
本当の深さの約4分の3です。水の中のストローのどの点も、同じ割合で浅く見えるので、ストローは水面で折れて、浅い向きに曲がって見えます。\\
極端な場合：$n=1$（水がない）なら $h'=h$ で、折れて見えません。
\kotae 約\qty{7.5}{cm}（本当の深さの約4分の3）
\end{phreidai}
\begin{jibun}
深さ\qty{30}{cm}の水そうの底の小石は、真上から見ると、何 cm の深さに見えますか。
\bkjans{$30\div1.333\approx\qty{22.5}{cm}$}
\end{jibun}

\label{ph:波-12:ex4}
\begin{phreidai}[ライフガードの最短時間の道（フェルマーの原理）]
砂浜のライフガード（水ぎわから\qty{20}{m}）が、海の人（水ぎわから\qty{20}{m}沖）を助けに行きます。2人は、水ぎわにそって\qty{30}{m}離れています。走る速さは\qty{6}{m/s}、泳ぐ速さは\qty{1.5}{m/s}です。水ぎわのどこから海に入れば、いちばん早く着きますか。そのとき、屈折の法則と同じ関係が成り立つことを確かめましょう。
\phbamen
光の屈折と同じ形の問題を、人の動きで考える場面です。
\phmushi
波で流されること、走り始めと泳ぎ始めの時間、砂の上の速さのむらを無視します。
\phzu
③の図\ref{ph:波-12:fig-fermat}。$h_1=h_2=\qty{20}{m}$、$L=\qty{30}{m}$、$v_1=6$、$v_2=\qty{1.5}{m/s}$。海に入る点を、ライフガードの真横から $x$〔m〕とする。
\phriyu
かかる時間 $T(x)$ がいちばん小さくなる $x$ を探します（フェルマーの原理、5-7）。表計算で $x$ を変えて調べられます。
\phtateru
\[T(x)=\frac{\sqrt{20^2+x^2}}{6}+\frac{\sqrt{20^2+(30-x)^2}}{1.5}\]
\phtoku
\begin{bkcenter}
\begin{bktable}{colspec={Q[l,m,wd=16mm] *{7}{Q[c,m,wd=10mm]}}, rowsep=1.2pt, colsep=4pt, row{1}={bg=bkTheme!10, font=\sffamily\small}}
$x$〔m〕 & 0 & 15 & 20 & 25 & 26 & 27 & 30 \\
$T$〔s〕 & 27.37 & 20.83 & 19.62 & 19.08 & 19.06 & 19.08 & 19.34 \\
\end{bktable}
\end{bkcenter}
$x\approx\qty{26}{m}$ でいちばん短く、約\qty{19.1}{s}です。そのとき、
\[\sin\theta_1=\frac{26}{\sqrt{20^2+26^2}}\approx0.793,\qquad\sin\theta_2=\frac{4}{\sqrt{20^2+4^2}}\approx0.196\]
\[\frac{\sin\theta_1}{\sin\theta_2}\approx4.0=\frac{v_1}{v_2}=\frac{6}{1.5}\]
\phtashikame
まっすぐ向かう道（$x=\qty{15}{m}$）は約\qty{20.8}{s}で、約1.8秒おそくなります。sin の比が速さの比に等しいという結果は、波-8 ①の屈折の法則と同じ形です。光は、この計算を毎回「自動で」していることになります。
\kotae 真横から約\qty{26}{m}の所で海に入る（約\qty{19.1}{s}）。$\sin\theta_1/\sin\theta_2\approx4=v_1/v_2$
\end{phreidai}
\begin{jibun}
例題4で、走る速さが\qty{5}{m/s}、泳ぐ速さが\qty{1}{m/s}なら、いちばん早い道で、$\sin\theta_1/\sin\theta_2$ はいくつになりますか。
\bkjans{$v_1/v_2=5\div1=5$}
\end{jibun}

\label{ph:波-12:ex5}
\begin{phreidai}[臨界角とダイヤモンドの輝き]
水（1.333）、ガラス（1.5）、ダイヤモンド（2.42）から空気へ出る光の臨界角を、それぞれ求めましょう。
\phbamen
宝石がよく光るわけを、全反射で考える場面です。
\phmushi
光の色による屈折率の違い（⑧）を無視します。
\phzu
④の図\ref{ph:波-12:fig-zenhansha}。出る側は空気（$n_2=1$）。
\phriyu
屈折角が\ang{90}になる入射角が臨界角（④）。
\phtateru
\[\sin\theta_{\mathrm c}=\frac1n\]
\phtoku
水：$\sin\theta_{\mathrm c}=1/1.333=0.750$、$\theta_{\mathrm c}\approx\ang{48.6}$\\
ガラス：$1/1.5=0.667$、$\theta_{\mathrm c}\approx\ang{41.8}$\\
ダイヤモンド：$1/2.42=0.413$、$\theta_{\mathrm c}\approx\ang{24.4}$
\phtashikame
屈折率が大きいほど臨界角は小さく、全反射が起きやすくなります。ダイヤモンドは、中に入った光の多くが全反射を繰り返し、上の面から出てくるので輝きます。
\kotae 水 約\ang{49}、ガラス 約\ang{42}、ダイヤモンド 約\ang{24}
\end{phreidai}
\begin{jibun}
氷（屈折率1.31）から空気へ出る光の臨界角を求めましょう。
\bkjans{$\sin\theta_{\mathrm c}=1/1.31=0.763$、約\ang{50}}
\end{jibun}

\label{ph:波-12:ex6}
\begin{phreidai}[光ファイバの受け入れ角]
コアの屈折率1.48、クラッドの屈折率1.46の光ファイバがあります。（1）コアとクラッドの境の臨界角、（2）開口数 NA と、はしから光を受け入れられる最大の角 $\theta_{\mathrm a}$ を求めましょう。
\phbamen
光通信のファイバに、光をどこまでななめに入れてよいかを決める場面です。
\phmushi
ファイバの曲がり、ガラスによる光の吸収、光が波であること（細いファイバでは効く）を無視します。
\phzu
④の図\ref{ph:波-12:fig-fiber}。$n_1=1.48$、$n_2=1.46$、外は空気。
\phriyu
境の臨界角は $\sin\theta_{\mathrm c}=n_2/n_1$、受け入れ角は $\sin\theta_{\mathrm a}=\sqrt{n_1^{\,2}-n_2^{\,2}}$（④）。
\phtateru
\[\sin\theta_{\mathrm c}=\frac{n_2}{n_1},\qquad\mathrm{NA}=\sin\theta_{\mathrm a}=\sqrt{n_1^{\,2}-n_2^{\,2}}\]
\phtoku
（1）$\sin\theta_{\mathrm c}=1.46\div1.48=0.9865$、$\theta_{\mathrm c}\approx\ang{80.6}$\\
（2）\[\mathrm{NA}=\sqrt{2.1904-2.1316}=\sqrt{0.0588}\approx0.242,\qquad\theta_{\mathrm a}=\sin^{-1}0.242\approx\ang{14.0}\]
\phtashikame
コアとクラッドの屈折率の差は約1.4\,\%だけですが、境の入射角が\ang{80.6}より大きい（境にほぼそって進む）光は全反射します。軸から\ang{14}以内の向きで入れれば、光は外へもれません。
\kotae （1）約\ang{80.6}、（2）NA 約0.24、約\ang{14}
\end{phreidai}
\begin{jibun}
コア1.50・クラッド1.47のファイバの NA と受け入れ角を求めましょう。
\bkjans{$\sqrt{2.25-2.1609}=\sqrt{0.0891}\approx0.30$、約\ang{17}}
\end{jibun}

\label{ph:波-12:ex7}
\begin{phreidai}[両凸レンズの焦点距離]
両面の半径がどちらも\qty{0.1}{m}の両凸レンズ（屈折率1.5）の焦点距離を求めましょう。平らな面と半径\qty{0.1}{m}の面からなる平凸レンズなら、どうなりますか。
\phbamen
ガラスをけずってレンズを作るとき、面の半径から焦点距離を決める場面です。
\phmushi
レンズの厚さ（薄いレンズとする）と、軸から遠い光の収差を無視します。
\phzu
⑥のレンズメーカーの式。光が左から来るとき、両凸では $R_1=\qty{0.1}{m}$、$R_2=\qty{-0.1}{m}$。平凸では平らな面の半径は限りなく大きい（$1/R=0$）。
\phriyu
レンズの面の半径と屈折率から焦点距離を出すのが、レンズメーカーの式（⑥）です。
\phtateru
\[\frac1f=(n-1)\Bigl(\frac1{R_1}-\frac1{R_2}\Bigr)\]
\phtoku
両凸：\[\frac1f=0.5\times\Bigl(\frac1{0.1}+\frac1{0.1}\Bigr)=0.5\times20=10,\qquad f=\qty{0.1}{m}\]
平凸：\[\frac1f=0.5\times\Bigl(\frac1{0.1}-0\Bigr)=5,\qquad f=\qty{0.2}{m}\]
\phtashikame
$n=1.5$ の両凸レンズは $f=R$ です（⑥）。平凸は、曲がった面が1つだけなので、曲げる力が半分、焦点距離が2倍です。
\kotae 両凸 \qty{0.1}{m}、平凸 \qty{0.2}{m}
\end{phreidai}
\begin{jibun}
両面の半径が\qty{0.1}{m}の両凸レンズで、屈折率が1.6なら、焦点距離はいくらですか。
\bkjans{$1/f=0.6\times20=12$、$f\approx\qty{0.083}{m}$}
\end{jibun}

\label{ph:波-12:ex8}
\begin{phreidai}[カメラのピント合わせ]
焦点距離\qty{50}{mm}のカメラのレンズで、\qty{2}{m}先に立つ身長\qty{1.7}{m}の人をうつします。（1）レンズから受光素子までの距離 $b$、（2）倍率と像の高さを求めましょう。（3）遠くの景色（$a$ が限りなく大きい）にピントが合っていたとき、レンズを何 mm 動かせばよいですか。
\phbamen
カメラがピントを合わせるとき、レンズをどれだけ動かすかを計算する場面です。
\phmushi
カメラのレンズを、焦点距離\qty{50}{mm}の1枚の薄いレンズと考えます。
\phzu
⑥の図\ref{ph:波-12:fig-sakuzu}。$f=\qty{0.05}{m}$、$a=\qty{2}{m}$。
\phriyu
物までの距離から像の距離を出すのは、レンズの式（⑥）。遠くの景色の像は焦点（$b=f$）にできます。
\phtateru
\[\frac1b=\frac1f-\frac1a,\qquad\text{倍率}=\frac ba\]
\phtoku
（1）\[\frac1b=\frac1{0.05}-\frac12=20-0.5=19.5,\qquad b=\frac{1}{19.5}\approx\qty{0.0513}{m}=\qty{51.3}{mm}\]
（2）倍率 $=0.0513\div2\approx0.026$、像の高さ $1.7\times0.0256\approx\qty{0.044}{m}=\qty{44}{mm}$（さかさま）\\
（3）$51.3-50=\qty{1.3}{mm}$ だけ、レンズを受光素子から遠ざける。
\phtashikame
物が遠いので、$b$ は焦点距離よりわずかに大きいだけです。人の像（\qty{44}{mm}）は、ふつうのカメラの受光素子（高さ\qty{24}{mm}など）からはみ出すので、全身をうつすにはもっと離れます。
\kotae （1）約\qty{51.3}{mm}、（2）約0.026倍、約\qty{44}{mm}、（3）約\qty{1.3}{mm}
\end{phreidai}
\begin{jibun}
例題8のカメラで、\qty{5}{m}先の物にピントを合わせるとき、$b$ と、遠くの景色から動かす量を求めましょう。
\bkjans{$1/b=20-0.2=19.8$、$b\approx\qty{50.5}{mm}$。約\qty{0.5}{mm}}
\end{jibun}

\label{ph:波-12:ex9}
\begin{phreidai}[2枚のレンズを行列で]
焦点距離\qty{0.1}{m}のレンズ2枚を、\qty{0.05}{m}はなして置きました。（1）全体の行列と合成焦点距離を求めましょう。（2）高さ\qty{0.01}{m}で軸に平行に入った光は、2枚めのレンズの何 m 後ろで軸と交わりますか。
\phbamen
カメラのように、レンズを組み合わせたときの働きを、1つの行列にまとめる場面です。
\phmushi
レンズの厚さと、軸から遠い光の収差を無視します。
\phzu
⑦の行列。すき間 $d=\qty{0.05}{m}$、$f_1=f_2=\qty{0.1}{m}$。光は1枚め→すき間→2枚めの順に通る。
\phriyu
順に通る要素の行列を、右から左へかけます（7-3・⑦）。
\phtateru
\[M=\begin{pmatrix}1&0\\-10&1\end{pmatrix}\begin{pmatrix}1&0.05\\0&1\end{pmatrix}\begin{pmatrix}1&0\\-10&1\end{pmatrix}\]
\phtoku
（1）右の2つから：\[\begin{pmatrix}1&0.05\\0&1\end{pmatrix}\begin{pmatrix}1&0\\-10&1\end{pmatrix}=\begin{pmatrix}0.5&0.05\\-10&1\end{pmatrix}\]
左からかけて：\[M=\begin{pmatrix}1&0\\-10&1\end{pmatrix}\begin{pmatrix}0.5&0.05\\-10&1\end{pmatrix}=\begin{pmatrix}0.5&0.05\\-15&0.5\end{pmatrix}\]
$C=-15$ なので、$f=1/15\approx\qty{0.067}{m}$。式 $1/f=10+10-0.05\times100=15$ とも合います。\\
（2）\[\begin{pmatrix}y\\\theta\end{pmatrix}=M\begin{pmatrix}0.01\\0\end{pmatrix}=\begin{pmatrix}0.005\\-0.15\end{pmatrix}\]
高さ0.005が0になるまで：$0.005\div0.15\approx\qty{0.033}{m}$。
\phtashikame
⑦の Python の結果と同じです。1枚のとき（$f=\qty{0.1}{m}$）より、2枚を組むと焦点距離が短く（曲げる力が強く）なりました。
\kotae （1）$M$ の成分は 0.5、0.05、$-15$、0.5、合成焦点距離 約\qty{0.067}{m}、（2）約\qty{0.033}{m}
\end{phreidai}
\begin{jibun}
焦点距離\qty{0.2}{m}と\qty{0.1}{m}のレンズを\qty{0.1}{m}はなして置いたときの合成焦点距離を求めましょう。
\bkjans{$1/f=5+10-0.1\div0.02=10$、$f=\qty{0.1}{m}$}
\end{jibun}

\label{ph:波-12:ex10}
\begin{phreidai}[虹の角]
赤い光（水の屈折率1.331）について、虹の光が集まる入射角 $\theta_1$ と、もどる角 $\phi=4\theta_2-2\theta_1$ を求めましょう。$\theta_1$ は、$\cos^2\theta_1=(n^2-1)/3$ で決まります。紫の光（1.343）では、もどる角は約\ang{40.6}です。虹の外側はどちらの色ですか。
\phbamen
雨上がりの虹の、赤と紫の位置を計算する場面です。
\phmushi
水の粒は完全な球とし、粒の表面で反射して返る光や、光の波としての性質（粒が小さいとき）を無視します。
\phzu
⑧の図\ref{ph:波-12:fig-niji}・図\ref{ph:波-12:fig-nijikaku}。$n=1.331$。
\phriyu
もどる角がいちばん大きくなる所に光が集まります。$\dd\phi/\dd\theta_1=0$（5-7）から求めた条件を使います。
\phtateru
\[\cos\theta_1=\sqrt{\frac{n^2-1}{3}},\qquad\sin\theta_2=\frac{\sin\theta_1}{n},\qquad\phi=4\theta_2-2\theta_1\]
\phtoku
\[\cos\theta_1=\sqrt{\frac{1.7716-1}{3}}=\sqrt{0.2572}\approx0.5071,\qquad\theta_1\approx\ang{59.5}\]
\[\sin\theta_2=\frac{0.8619}{1.331}\approx0.6476,\qquad\theta_2\approx\ang{40.4}\]
\[\phi=4\times40.36-2\times59.53\approx\ang{42.4}\]
\phtashikame
赤 \ang{42.4} は紫 \ang{40.6} より大きいので、赤のほうが、頭の影の向きから遠い所（虹の外側）に見えます。Python で $\theta_1$ を少しずつ変えて $\phi$ のいちばん大きい値を探しても、同じ\ang{42.4}になります。
\kotae 入射角 約\ang{59.5}、もどる角 約\ang{42.4}。外側が赤
\end{phreidai}
\begin{jibun}
黄色の光（$n=1.333$）では、虹の光が集まる入射角ともどる角はいくらですか。
\bkjans{$\cos\theta_1=\sqrt{0.2590}\approx0.509$、$\theta_1\approx\ang{59.4}$、$\theta_2\approx\ang{40.2}$、$\phi\approx\ang{42.1}$}
\end{jibun}

\label{ph:波-12:ex11}
\begin{phreidai}[虫めがねの倍率]
焦点距離\qty{10}{cm}の虫めがねの倍率を求めましょう。（1）物を焦点の所に置いて、目を楽にして見るとき、（2）像が明視の距離（\qty{25}{cm}）にできるように、虫めがねを目に近づけて見るとき。
\phbamen
小さな字や虫を、虫めがねで大きく見る場面です。
\phmushi
虫めがねと目の間の距離（目にくっつけるとする）、レンズの収差を無視します。
\phzu
⑥の図\ref{ph:波-12:fig-kyozou}・⑨。$f=\qty{10}{cm}$、明視の距離\qty{25}{cm}。
\phriyu
倍率は、目で見る角が、明視の距離で見るときの何倍になるかです（⑨）。
\phtateru
\[（1）\ \frac{25}{f}\qquad（2）\ \frac{25}{f}+1\]
\phtoku
（1）$25\div10=2.5$ 倍。（2）$2.5+1=3.5$ 倍。\\
（2）を確かめる：像を $b=\qty{-25}{cm}$ にするには、$1/a=1/10+1/25=0.14$ より $a\approx\qty{7.14}{cm}$。倍率 $|b|/a=25\div7.14=3.5$。
\phtashikame
焦点距離が短いほど、倍率が大きくなります。$f=\qty{5}{cm}$ なら5倍（6倍）です。
\kotae （1）2.5倍、（2）3.5倍
\end{phreidai}
\begin{jibun}
焦点距離\qty{5}{cm}の虫めがねの、例題11の（1）（2）の倍率を求めましょう。
\bkjans{（1）$25\div5=5$ 倍、（2）6倍}
\end{jibun}

\label{ph:波-12:ex12}
\begin{phreidai}[望遠鏡の倍率と長さ]
焦点距離\qty{900}{mm}の対物レンズと、焦点距離\qty{20}{mm}の接眼レンズの望遠鏡があります。倍率と、筒のおよその長さを求めましょう。接眼レンズを焦点距離\qty{10}{mm}のものにかえると、どうなりますか。
\phbamen
天体望遠鏡で月や木星を見るとき、接眼レンズを選ぶ場面です。
\phmushi
レンズの厚さ、収差、空気のゆらぎを無視します。
\phzu
⑨の図\ref{ph:波-12:fig-bouenkyou}。$f_{\text{対物}}=\qty{900}{mm}$、$f_{\text{接眼}}=\qty{20}{mm}$。
\phriyu
望遠鏡の倍率は $f_{\text{対物}}/f_{\text{接眼}}$、筒の長さは2つの焦点距離の和（⑨）。
\phtateru
\[\text{倍率}=\frac{f_{\text{対物}}}{f_{\text{接眼}}},\qquad\text{長さ}\approx f_{\text{対物}}+f_{\text{接眼}}\]
\phtoku
\[\frac{900}{20}=45\ \text{倍},\qquad900+20=\qty{920}{mm}\]
接眼\qty{10}{mm}なら $900\div10=90$ 倍、長さ約\qty{910}{mm}。
\phtashikame
接眼レンズの焦点距離を半分にすると、倍率は2倍です。ただし、倍率を上げると、見える範囲がせまく、暗くなります。
\kotae 45倍、約\qty{92}{cm}。接眼\qty{10}{mm}で90倍
\end{phreidai}
\begin{jibun}
対物レンズの焦点距離\qty{700}{mm}、接眼レンズ\qty{10}{mm}の望遠鏡の倍率と、筒のおよその長さを求めましょう。
\bkjans{70倍、約\qty{710}{mm}}
\end{jibun}

\label{ph:波-12:ex13}
\begin{phreidai}[机の上の照度]
光度\qty{1000}{cd}の小さな照明が、机の上\qty{2}{m}の高さにあります。（1）真下の机の照度を求めましょう。（2）真下から横に\qty{1.5}{m}離れた所の、机の面の照度を求めましょう。
\phbamen
照明の明るさが、机の場所によってどう変わるかを見積もる場面です。
\phmushi
照明を1点とし、どの向きにも同じ光度で光を出すとします。壁や天井で反射して届く光を無視します。
\phzu
⑩。$I=\qty{1000}{cd}$。（2）では、照明までの距離 $r=\sqrt{2^2+1.5^2}=\qty{2.5}{m}$、机の法線と光のなす角 $\theta$ は $\cos\theta=2/2.5=0.8$。
\phriyu
照度の逆2乗の法則と、面の傾きの $\cos\theta$（⑩）。
\phtateru
\[E=\frac{I\cos\theta}{r^2}\]
\phtoku
（1）\[E=\frac{1000\times1}{2^2}=\qty{250}{lx}\]
（2）\[E=\frac{1000\times0.8}{2.5^2}=\frac{800}{6.25}=\qty{128}{lx}\]
\phtashikame
（2）は（1）の約半分です。距離が1.25倍で $1/1.25^2=0.64$、傾きで0.8、あわせて $0.64\times0.8\approx0.51$ 倍。事務室の目安\qty{750}{lx}には足りないので、照明を増やすか、近づけます。\\
単位：$\text{cd}/\text{m}^2=\text{lm}/(\text{sr}\cdot\text{m}^2)$ で、立体角 sr は単位のない量なので lx になります。
\kotae （1）\qty{250}{lx}、（2）\qty{128}{lx}
\end{phreidai}
\begin{jibun}
光度\qty{400}{cd}の照明から\qty{2}{m}離れ、光に垂直な面の照度を求めましょう。距離を\qty{1}{m}にすると、照度は何倍になりますか。
\bkjans{$400\div4=\qty{100}{lx}$。距離が半分で4倍の\qty{400}{lx}}
\end{jibun}

%% ======================== まちがえやすいところ ========================
\Needspace{11\baselineskip}\bksection{tsumazuki}{まちがえやすいところをチェック}
\begin{tsumazuki}{入射角は、どこからはかる？}
\bkng{入射角は、光と境の面（水面・鏡の面）のなす角。}
\bkok{入射角・反射角・屈折角は、どれも\textbf{法線}（面に垂直な線）からはかる。}
面からはかった\ang{30}は、法線からは\ang{60}です。図に法線を引いてから角を書きましょう（②③）。
\end{tsumazuki}

\begin{tsumazuki}{遅い物に入ると、どちらへ曲がる？}
\bkng{水に入ると、光は水面にそう向きに曲がる。}
\bkok{遅い物（屈折率が大きい物）に入ると、法線に\textbf{近づく}。速い物へ出ると、法線から遠ざかる。}
ライフガードは、遅い海の中の道を短くしようとします（③）。
\end{tsumazuki}

\begin{tsumazuki}{全反射は、どちらから出るとき？}
\bkng{空気から水へ入る光でも、角が大きいと全反射する。}
\bkok{全反射は、屈折率の\textbf{大きい}物から\textbf{小さい}物へ出るとき（水から空気へ）だけ起きる。}
空気から水へは、どんな角でも入れます（屈折角は臨界角以下になる。④）。
\end{tsumazuki}

\begin{tsumazuki}{レンズの式の $b$ が負}
\bkng{$b$ が負になったので、計算をまちがえた。}
\bkok{$b<0$ は、物と同じ側に虚像ができるという意味。凹レンズは $f$ を負にして同じ式を使う。}
\end{tsumazuki}

\begin{tsumazuki}{照度は距離に反比例？}
\bkng{距離が2倍なら、照度は半分。}
\bkok{照度は距離の\textbf{2乗}に反比例するので、$\tfrac14$。}
光束が、半径の2乗に比例する球の面に広がるからです（⑩）。
\end{tsumazuki}

\begin{column}{フェルマーと「いちばん短い時間」}
ピエール・ド・フェルマー（1607〜1665年）は、フランスの法律家で、数学を趣味として研究した人です。数学では「フェルマーの最終定理」で知られています。1662年ごろ、フェルマーは、光が通る道は「かかる時間がいちばん短い道」だと考え、屈折の法則を導きました。この考えだと、光は水やガラスの中で、空気の中より遅くなければなりません。光が水の中で空気の中より遅くなることを実験で確かめたのは、およそ190年後の1850年、フーコーでした（回る鏡を使った実験）。のちに、「何かがいちばん小さく（または大きく）なる道を、自然は選ぶ」という考えは、力学のラグランジュの式（波-3）や、量子力学にまで広がりました。
\end{column}

%% ======================== 工学ではこう使う ========================
\Needspace{10\baselineskip}\bksection{phgear}{工学ではこう使う}
\begin{phkougaku}[光ファイバ通信・スマホのカメラ・めがね]
\textbf{光ファイバ通信}：インターネットのデータの多くは、海底や地下の光ファイバの中を、光の点めつとして運ばれています。全反射（④）で光を閉じこめるので、何十 km 進んでも光が外へもれません。ただし、角の大きい光は、軸にそう光より長い道を進むので、同じ時刻に出た光がばらばらの時刻に着きます（コア1.48・クラッド1.46なら、\qty{1}{km}で約\qty{68}{ns}の差）。長い距離では、コアを数 \si{\micro m} と細くして、光が1つの道だけを進むようにしたファイバ（シングルモード）を使います。

\textbf{スマホのカメラ}：スマホのレンズは、数 mm の厚さに、何枚ものプラスチックのレンズを重ねたものです。収差を打ち消すように1枚ずつの形を決める設計は、⑦の光線の追いかけを、コンピュータで何万本も繰り返して行います。

\textbf{めがね}：めがね店では、レンズの度数（⑨、D）で注文します。$-2$\,D のレンズは焦点距離\qty{-0.5}{m}の凹レンズで、\qty{0.5}{m}より遠くがぼやける近視の目に合います。

\begin{phketa}[光の数の目安]
\phketaline{屈折率}{空気 1.0003、水 1.33、ガラス 約1.5、ダイヤモンド 2.42}
\phketaline{臨界角}{水 \ang{49}、ガラス \ang{42}、ダイヤモンド \ang{24}}
\phketaline{焦点距離}{スマホのレンズ 数 mm、カメラ 50 mm、虫めがね 10 cm}
\phketaline{照度}{事務室 750 lx、晴れた日の直射日光 約10万 lx}
\phboxhead{答えが桁で合っているか}
屈折率は1以上。$\sin\theta>1$ なら全反射。照度では $r^2$ を忘れずに。
\end{phketa}
\end{phkougaku}
\phstation{光学機器,光通信工学,照明工学,画像工学}{屈折率、臨界角、開口数（NA）、焦点距離、収差、光線行列、光束・光度・照度}



%% ======================== 章末問題 ========================
\bkendproblems
\bkproblemgroup{基本}{3}
\bkproblem 空気からガラス（屈折率1.5）へ、入射角\ang{60}で入る光の屈折角を求めましょう。
\bkpans{$\sin\theta_2=0.866\div1.5=0.577$、約\ang{35}}
\bkproblem 深さ\qty{1.2}{m}のプールの底は、真上から見ると何 m の深さに見えますか。
\bkpans{$1.2\div1.333\approx\qty{0.90}{m}$}
\bkproblem 屈折率1.65のガラスから空気へ出る光の臨界角を求めましょう。
\bkpans{$\sin\theta_{\mathrm c}=1/1.65=0.606$、約\ang{37}}
\bkproblem 焦点距離\qty{0.15}{m}の凸レンズの前、（1）\qty{0.45}{m}、（2）\qty{0.10}{m}の所に物を置きます。像の位置と倍率、実像か虚像かを答えましょう。
\bkpans{（1）$1/b=1/0.15-1/0.45$ より $b=\qty{0.225}{m}$、倍率0.5、さかさまの実像。（2）$b=\qty{-0.30}{m}$（レンズの手前）、倍率3、正立の虚像}
\bkproblem 身長\qty{150}{cm}の人が全身を映すには、鏡の長さは少なくとも何 cm いりますか。
\bkpans{$150\div2=\qty{75}{cm}$}
\bkproblem 光度\qty{200}{cd}の小さな照明から、光に垂直な面までの距離が（1）\qty{1.0}{m}、（2）\qty{2.0}{m}のときの照度を求めましょう。
\bkpans{（1）$200\div1^2=\qty{200}{lx}$、（2）$200\div2^2=\qty{50}{lx}$}
\bkproblemgroup{チャレンジ}{4}
\bkproblem 平凸レンズ（平らな面と半径\qty{0.05}{m}の面、屈折率1.5）の焦点距離は？ 水（1.333）の中では約何倍ですか（⑥の「レンズの屈折率と周りの屈折率の差を2倍にすると？」）。
\bkpans{空気中：$1/f=0.5\times20=10$、$f=\qty{0.10}{m}$。水中：$1/f=(1.5/1.333-1)\times20\approx2.5$、$f\approx\qty{0.40}{m}$、約4倍}
\bkproblem 焦点距離\qty{0.05}{m}の凸レンズの\qty{0.02}{m}後ろに、焦点距離 $\qty{-0.1}{m}$ の凹レンズを置きました。合成焦点距離を求めましょう。
\bkpans{$1/f=20-10-0.02\div(0.05\times(-0.1))=10+4=14$、$f\approx\qty{0.071}{m}$}
\bkproblem 焦点距離\qty{1000}{mm}の対物レンズの望遠鏡で、40倍にしたい。接眼レンズの焦点距離と、筒のおよその長さを求めましょう。
\bkpans{$1000\div40=\qty{25}{mm}$、長さ約\qty{1025}{mm}}
\bkproblem 光ファイバのコアの屈折率が1.50のとき、受け入れ角を\ang{10}以上にするには、クラッドの屈折率をいくつ以下にすればよいですか（$\sin\ang{10}\approx0.174$）。
\bkpans{$1.50^2-n_2^{\,2}\geqq0.174^2$、$n_2\leqq\sqrt{2.25-0.0302}\approx1.4899$。1.48以下（上限なので切り捨て）}

\begin{furikaeri}
\begin{checklist}
\item 光の直進・反射・屈折を、身近な例と作図で説明できる。\bkfurref{ph:波-12:sec1}
\item 屈折の法則で、屈折角・臨界角・見かけの深さを求められる。\bkfurref{ph:波-12:sec3}
\item レンズの式で、像の位置と倍率を求め、光学機器の倍率を見積もれる。\bkfurref{ph:波-12:sec6}
\item 屈折率が色で変わること（分散）から、プリズムと虹の色を説明できる。\bkfurref{ph:波-12:sec8}
\item フェルマーの原理から反射・屈折の法則を導き、光線の行列でレンズの組を計算できる。\bkfurref{ph:波-12:sec7}
\end{checklist}
\end{furikaeri}
\bkfukidashi{mina}{ストローが曲がって見えたのは、光が「早く着く道」を選んでいたからなんだ。虹の\ang{42}まで計算できるなんて、おどろいた！}
